% Modernised reading — fonds Alexandre Grothendieck, shelfmark 101
% (pages 1-42, the whole folder).
%
% This is an INTERPRETATION, not a transcription. It re-reads the pages in
% today's notation and vocabulary, states the hypotheses the manuscript leaves
% implicit, and drops the critical apparatus so that the mathematics reads
% straight through. Everything the brackets and underlines carried moves into
% a footnote. It is held to being mathematically correct as it stands; where
% the manuscript is loose or mistaken, the true statement is what appears here
% and the footnote says what the page has. In any dispute of substance the
% transcription governs: whoever cites Grothendieck cites batch-01.fr.tex
% (pp. 1-20), batch-02.fr.tex (pp. 21-40) or batch-03.fr.tex (pp. 41-42).
%
% Pass: Opus 5.5 (claude-opus-5-5), 2026-10-10 — first pass, unchecked against the pages by a human.
% The folder was read whole, its three batches together; this is one document
% for the shelfmark. It was made on Opus 5.5 and has not been measured against
% the reference reading of folder 115 (made on Opus 5). The literature named
% in the footnotes is cited from memory and was not checked against the
% sources.
%
% Spine. Three runs, each opened by its own buff cover (pp. 1, 5, 16):
%   I    pp. 2-4    local cohomology on the punctured cone vs. graded Ext on
%                   Proj; fragmentary.
%   II   pp. 6-15   « Dualité locale pour Modules non quasi cohérents »: Ext
%                   on Spf A (6-7), scattered notes (8), local duality for
%                   arbitrary O_X-modules on Spec of a complete regular local
%                   ring: true for n <= 1, fails for n = 2 with algebraic
%                   duals (the page's own conclusion, 14-15).
%   III  pp. 17-34  « Notes antiques 1958 »: plan (17), module fondamental
%                   (= canonical module) and local duality (18-21), the
%                   pairing and Prop. 9-12 (22-29), « 6. » Gorenstein rings
%                   (30-34).
%   IV   pp. 35-39  first recapitulation, Theorem, Problems A, B, C.
%   V    pp. 40-42  second recapitulation, exactness of Ext (41), the case
%                   M = A (42). IV and V sit under the same cover as III.
% That pp. 21-34 are articles 5 and 6 of § 2 of the plan of p. 17 is THIS
% READING's observation (p. 30's heading « 6. ... Anneaux de Gorenstein »
% matches § 2, n° 6 of the plan); the leaves do not say so.
%
% Notation fixed for the folder. (A, m, k) local noetherian, n = dim A;
% H^i(M) = H^i_m(M) as on the pages; I = E(k), his « module dualisant »;
% D = Hom_A(-, I); « foncteur dualisant » = isomorphic to D on modules of
% finite length; Ω_A his « module fondamental » = today's canonical module.
% His gothic N names three different ideals; here: 𝔑 the annihilator of Ω
% (Prop. 9), 𝔞 the kernel of A' -> A (pp. 26-27), 𝔮 an ideal of definition
% (pp. 35-42). His « corang n » (uncertain reading) is rendered « de
% dimension n » (dim A/p = n). His « codim » of a module is the depth. In
% section I the twist index is j (the page's n).
%
% Substantive departures from the pages, with pages.
%  - p. 2: noetherian hypothesis supplied for Ext to commute with the sum.
%  - p. 10: the claim « isomorphismes pour tout i » is false for n >= 2 with
%    algebraic duals; pp. 14-15 find this themselves. Stated as such.
%  - p. 21 (Cor. 5) and p. 28 (Lemme): inequalities on depth corrected
%    (codim M > k, resp. >= k; Ext vanishing for i != n - k, not i != k).
%  - p. 27: (2), (2 bis), (13) have M where A is meant; « dualisant » where
%    « fondamental » is meant. p. 28, Prop. 11: localisation of Ω holds for p
%    in Supp Ω (or A equidimensional), not for every p; counterexample given.
%  - pp. 30-36: the equivalences are stated under Ext^i(k,M) = 0 for all
%    i < n (depth M >= n), which the proofs use; the pages assume only
%    Ext^{n-1}(k,M) = 0. Later literature footnoted.
%  - p. 32, Prop. 15 (d): « équidimensionnel » must be « all associated
%    primes of  of dimension n »; counterexample k[[x,y]]/(x^2,xy).
%  - p. 33: struck Th. 5 (« Gorenstein » in the sense of Prop. 15 ⇒ CM) is
%    false (cone over an abelian surface); the folder's last page (42)
%    attempts the same implication from (vii), and its proof stops on the
%    exact sequence whose third term is the obstruction. Stated as such.
%  - p. 36, Problem A: false as stated (k[[x,y]]/(x^2,xy)); p. 37 Problem C
%    and p. 40 « Pb. »: true for M of finite type, by results of 1963-1987.
%  - p. 39: H^n(A/xA) read H^{n-1}(A/xA).
%  - pp. 40-42: the hat read as the dual D (by comparison with p. 35 (iv)),
%    not completion as on p. 22. The bar of p. 42 is undefined; (vii) read
%    as A ≅ Ω_A.
%  - p. 41: the « ou bien » alternative is unnecessary; the equivalence
%    Ext^i(k,M) = 0 <=> H^i(M) = 0 holds under the hypothesis alone.
%
% Modern names given and footnoted as ours: Matlis duality, local duality,
% canonical module, deficiency modules (Schenzel), S2-ification,
% quasi-Gorenstein ring, Bass numbers, Serre-Grothendieck correspondence,
% Rees's lemma. « Anneau de Gorenstein » is on the page (p. 34).
%
% What the folder announces and never establishes: Th. 1, Th. 3 and § 1
% (cited, not in the folder); Th. 4's statement (p. 18, illegible); the
% sequel of p. 4 (« Cas particuliers »), of p. 7, and of p. 15 (topological
% duals); Prop. 14's induction (p. 31); Prop. 17's Cor. (p. 34, no verb);
% Problems A, B, C (pp. 36-37); the step of p. 42.
%
% Group. The folder sits in the inventory group « Géométrie et topologie
% combinatoire » (69-102); its subject does not match (as for folder 100).
% It meets none of the modernised folders of the group (81, 88, 89).
\documentclass[11pt,a4paper]{article}
\input{../preamble/grothendieck}

\folder{101}
\pages{1}{42}
\dating{[à partir de 1967-1968]}
\watermark{Édition de démonstration\\interprétation personnelle de l'œuvre}
\foldertitle{Cohomologie locale : notes manuscrites (s.d.) — lecture modernisée du dossier entier}

\begin{document}

\section*{Résumé}

\begin{resume}
Ces feuillets parlent de ce qui se passe \emph{près d'un point} : comment un
objet algébrique --- un anneau, un module --- se comporte au voisinage d'un
point fermé, une fois le point lui-même retiré. L'instrument est la
\emph{cohomologie locale}, qui mesure ce qu'on perd ou ce qu'on gagne en
retirant le point, un peu comme la topologie compare une boule pleine à la
même boule trouée.

Le problème est un problème de \emph{dualité}. En algèbre linéaire, un espace
de dimension finie et son dual se déterminent l'un l'autre, et un accouplement
non dégénéré les met face à face. Grothendieck cherche l'énoncé analogue pour
un anneau local : la cohomologie locale d'un module, en degré $p$, devrait être
le « dual » d'un certain groupe $\mathrm{Ext}$ en degré complémentaire $n-p$,
où $n$ est la dimension. Pour que cela marche, il faut un module de référence
contre lequel dualiser, comme il faut le corps de base pour dualiser un espace
vectoriel. Il l'appelle \emph{module fondamental} ; on l'appelle aujourd'hui
\emph{module canonique}. La dualité tient alors exactement quand l'anneau est
de Cohen-Macaulay --- quand sa dimension est « honnête » partout, sans parties
de dimension plus petite qui traînent ou qui s'accrochent. Et parmi ces
anneaux, ceux qui sont leur propre module canonique, comme un espace muni d'un
produit scalaire est son propre dual, sont les \emph{anneaux de Gorenstein}.

Le dossier contient trois textes. Le premier, très bref, compare la cohomologie
d'une variété projective à la cohomologie locale du cône au-dessus d'elle. Le
deuxième essaie d'étendre la dualité locale à \emph{tous} les faisceaux de
modules, même ceux qui ne viennent pas d'un module ; l'essai réussit en
dimension $1$ et échoue en dimension $2$, et Grothendieck le constate : il
faudrait des duaux topologiques. Le troisième, le plus long, porte sur sa
chemise « Notes antiques 1958 » : un plan, la définition du module fondamental,
le théorème de dualité locale, puis une longue chasse aux conditions qui
caractérisent le module canonique d'un anneau de Cohen-Macaulay et les anneaux
de Gorenstein. On y voit une méthode : il dresse la liste des conditions,
dessine le graphe des implications, raye, recommence, isole ce qui reste sous
le nom de « Problème A, B, C ». Plusieurs de ces problèmes ont été résolus
bien plus tard ; l'un a une réponse négative. Le dossier s'arrête au milieu
d'une démonstration qui ne peut pas aboutir, et sur la ligne même où se
trouve l'obstruction.

Pour aller plus loin : dualité de Matlis, dualité locale, module canonique,
anneaux de Cohen-Macaulay et de Gorenstein, modules de déficience, nombres de
Bass.

\keywords{local cohomology, Matlis duality, local duality, canonical module, Cohen-Macaulay ring, Gorenstein ring, quasi-Gorenstein ring, deficiency modules, Bass numbers, S2-ification, Serre-Grothendieck correspondence, formal scheme, local duality for non-quasi-coherent sheaves}
\end{resume}

\section*{Le fil du dossier, et les conventions}

\subsection*{Les stations}

\begin{enumerate}
\item[I.] \emph{Cohomologie locale et cohomologie projective, via les cônes}
(pages 2 à 4 ; chemise page 1) : $\mathrm{Ext}$ à supports sur le cône épointé
au-dessus de $X = \operatorname{Proj}(\mathcal{S})$ et sommes graduées
d'$\mathrm{Ext}$ sur $X$.
\item[II.] \emph{Dualité locale pour Modules non quasi cohérents} (pages 6 à
15 ; chemise page 5) : les $\mathrm{Ext}$ sur un schéma formel affine (pages
6 et 7), des notes éparses (page 8), puis la dualité locale pour un
$\mathcal{O}_X$-Module quelconque sur le spectre d'un anneau local régulier
complet : vraie en dimension $\leqslant 1$, fausse en dimension $2$ avec les
duaux algébriques (pages 10 à 15).
\item[III.] \emph{Notes antiques 1958} (pages 17 à 34 ; chemise page 16) :
\end{enumerate}
\begin{enumerate}
\item[III.1.] le plan, le module fondamental, le théorème de dualité locale et
ses corollaires (pages 17 à 21) ;
\item[III.2.] l'accouplement $\mathfrak{D}$, l'annulateur de $\Omega$, les
anneaux quotients d'un anneau de Cohen-Macaulay, les modules $\Omega_A^{(i)}$
et une suite spectrale (pages 22 à 29) ;
\item[III.3.] « 6. Caractérisation des modules fondamentaux d'anneaux Coh.\
Mac. Anneaux de Gorenstein » (pages 30 à 34).
\end{enumerate}
\begin{enumerate}
\item[IV.] \emph{Récapitulation, théorème, problèmes} (pages 35 à 39), sous
la même chemise.
\item[V.] \emph{Seconde récapitulation, exactitude des foncteurs
$\mathrm{Ext}$, et le cas $M = A$}, où le dossier s'arrête (pages 40 à 42).
\end{enumerate}

Les trois textes ne se citent pas. Le premier et le deuxième sont des
brouillons de recherche ; le troisième est, jusqu'à la page 34, la rédaction
suivie d'un chapitre, que deux récapitulations prolongent (IV et V). Les pages 9 et 13 sont des feuillets
dactylographiés d'un autre texte, sans trace de sa main, et ne sont pas lues ;
les pages 8 et 12 sont écrites au verso de tapuscrits, et seules ses notes
sont lues.

Le sujet du dossier n'a rien de commun avec celui du groupe d'inventaire où il
est classé, « Géométrie et topologie combinatoire », ni avec les dossiers 81,
88 et 89 de ce groupe, qui parlent de cartes, de dessins et d'ensembles
finis ; aucun renvoi n'est donc fait vers eux.

\subsection*{Ce que seule une lecture d'ensemble peut dire}

Le plan de la page 17 annonce, pour son § 2, six articles, dont le cinquième
est « Module fondamental et th.\ de dualité locale » et le sixième « Anneaux de
Gorenstein ». Le titre que porte la page 30, « 6. Caractérisation des modules
fondamentaux d'anneaux Coh.\ Mac.\ Anneaux de Gorenstein », est celui du
sixième article ; les pages 18 à 29, qui définissent le module fondamental et
en tirent la dualité, sont donc vraisemblablement le cinquième. C'est une
observation de cette lecture : les feuillets ne le disent pas, et les
théorèmes 1 et 3 ainsi que le § 1 qu'ils citent ne sont pas dans le dossier.

L'énoncé sur lequel le dossier s'arrête, page 42 --- qu'un anneau qui est son
propre module fondamental est de Cohen-Macaulay ---, est celui du « Théorème
5 » biffé de la page 33. Entre les deux, la page 34 a pris le parti inverse et
\emph{défini} un anneau de Gorenstein comme un anneau de Cohen-Macaulay qui est
son propre module fondamental. L'énoncé est faux en général
(page 42 ci-dessous) ; la définition de la page 34 est celle qu'on a gardée.

Les « Problèmes » des pages 36 et 37 reviennent page 40 sous le nom « Pb. »,
et la question de la page 41 en est une variante.

\subsection*{Conventions, valables pour tout le dossier}

$A$ est un anneau local noethérien, $\mathfrak{m}$ son idéal maximal,
$k = A/\mathfrak{m}$, $n = \dim A$, et $\widehat{A}$ son complété. Pour un
$A$-module $M$, $H^i(M) = H^i_{\mathfrak{m}}(M)$ est la cohomologie locale à
supports dans le point fermé, comme sur les pages. $I = E(k)$ est l'enveloppe
injective de $k$ --- le « module dualisant » des pages ---, et
$D(M) = \operatorname{Hom}_A(M, I)$ le dual de Matlis\footnote{La dualité de
Matlis est publiée en 1958 (E.~Matlis, \emph{Injective modules over Noetherian
rings}) ; les pages ne citent personne. Elles disent « module dualisant » pour
$I$ et « $H^n(M)$ est dualisant » pour $H^n(M) \simeq I$ ; le mot désigne
aujourd'hui tout autre chose (le module canonique d'un anneau de
Cohen-Macaulay), d'où le choix de ne pas le reprendre.}. Pour $M$ de type fini,
$D(D(M)) \simeq \widehat{M}$, et $D$ est une anti-équivalence des modules de
longueur finie sur eux-mêmes. Un foncteur contravariant des $A$-modules de
longueur finie vers les $A$-modules est dit \emph{dualisant}, comme sur les
pages, s'il est isomorphe à $D$. Un module est \emph{comonogène} si son socle
$\operatorname{Hom}(k, M)$ est de dimension $1$ et $M$ est de
$\mathfrak{m}$-torsion, c'est-à-dire s'il se plonge dans $I$ comme sous-module
essentiel.

$\Omega_A$, le « module fondamental » de la page 18, est un $A$-module de type
fini tel que $D(\Omega_A) \simeq H^n(A)$ ; c'est le \emph{module
canonique}\footnote{Pour un anneau de Cohen-Macaulay, J.~Herzog et E.~Kunz,
\emph{Der kanonische Modul eines Cohen-Macaulay-Rings} (1971) ; pour un anneau
local quelconque, avec exactement cette définition ($\Omega \otimes \widehat{A}
\simeq D(H^n(A))$), Y.~Aoyama (1980, 1983). Le nom « module fondamental » ne
s'est pas conservé. Cité de mémoire.} d'aujourd'hui, souvent noté $\omega_A$ ou
$K_A$. On garde $\Omega_A$, la lettre des pages.

L'idéal que les pages notent d'une même majuscule gothique désigne trois
idéaux différents : on écrit ici $\mathfrak{N}$ pour l'annulateur de $\Omega$
(pages 22 à 25), $\mathfrak{a}$ pour le noyau d'une surjection $A' \to A$
(pages 26 et 27), $\mathfrak{q}$ pour un idéal de définition (pages 35 à 42).
Un idéal premier « de corang $n$ »\footnote{« corang » est une lecture
incertaine d'un mot bref devant « $n$ », tenue la même dans tout le dossier ;
le sens, lui, n'est pas douteux : la page 38 écrit « corang » au-dessus de
« dimension » biffé.} est ici un premier $\mathfrak{p}$ \emph{de dimension
$n$}, c'est-à-dire tel que $\dim A/\mathfrak{p} = n$. La « codim » d'un module,
et la « codim.\ coh. » de la page 41, sont sa \emph{profondeur}
$\operatorname{prof} M$, le plus petit $i$ tel que $H^i(M) \neq 0$ ; $M$ de
type fini est de Cohen-Macaulay de dimension $n$ si $\operatorname{prof} M =
\dim M = n$. $\mathrm{lg}_{\mathfrak{p}}(M)$ est la longueur de
$M_{\mathfrak{p}}$ sur $A_{\mathfrak{p}}$, pour $\mathfrak{p}$ minimal.

Lorsque $\Omega_A$ intervient, on suppose qu'il existe, ce qui est le cas si
$A$ est complet ou quotient d'un anneau régulier (page 28), et, quand une
démonstration le demande, que $A$ est complet --- « ce qui est loisible »,
dit la page 36.

\subsection*{Ce que le dossier annonce sans l'établir}

Les théorèmes 1 et 3, le § 1 et sa « prop.\ 4, cor.\ 1 » sont cités et ne sont
pas dans le dossier. L'énoncé du théorème 4 (page 18) ne se lit pas ; ce qu'on
en sait vient de ses corollaires (page 20). Ne sont pas menés à terme : les
« cas particuliers » de la page 4, le calcul de la page 7, la dualité
topologique annoncée page 15, la récurrence de la proposition 14 (page 31),
le corollaire sans verbe de la proposition 17 (page 34), les Problèmes A, B, C
(pages 36 et 37), la démonstration de la page 42.

\pagerange{2}{4}
\section*{I. Cohomologie locale et cohomologie projective, via les cônes (pages 2 à 4)}

Soient $Y$ un préschéma noethérien, $\mathcal{S}$ une $\mathcal{O}_Y$-Algèbre
graduée quasi-cohérente à degrés $\geqslant 0$, engendrée par $\mathcal{S}_1$,
qui est de type fini\footnote{L'hypothèse noethérienne est ajoutée : elle sert
plus bas à faire commuter les $\mathrm{Ext}$ aux sommes directes. La page
écrit « préschéma de base ».}, $X = \operatorname{Proj}(\mathcal{S})$,
$C = \operatorname{Spec}(\mathcal{S})$ le cône affine et $X'$ le cône épointé,
complémentaire du sommet. La projection $p : X' \to X$ est affine et plate, et
\[
  X' \simeq \mathbf{V}(\mathcal{O}_X(1)) - \sigma(X)
  = \operatorname{Spec}_X\Bigl(\bigoplus_{j \in \mathbf{Z}} \mathcal{O}_X(j)\Bigr),
\]
où $\sigma$ est la section nulle\footnote{Avec la convention d'EGA,
$\mathbf{V}(\mathcal{E}) = \operatorname{Spec} \operatorname{Sym}(\mathcal{E})$.
La page écrit $\coprod_n \mathcal{O}(-n)$, qui est la même somme ; l'indice de
torsion, $n$ sur la page, est noté $j$ dans cette section. Le mot biffé après
« Spec » est illisible.}.

Soient $Z$ un fermé et $U$ un ouvert de $X$, $T = Z \cap U$ (localement fermé
dans $X$), $U' = p^{-1}(U)$, $T' = p^{-1}(T)$, et $F$, $G$ deux
$\mathcal{O}_U$-Modules quasi-cohérents, $F$ de type fini. Alors
\[
  \boxed{\ \operatorname{Ext}^i_{T'}(U'; p^{*}F, p^{*}G) \simeq
  \bigoplus_{j \in \mathbf{Z}} \operatorname{Ext}^i_{T}(U; F, G(j))\ }
\]
où $\operatorname{Ext}_T$ désigne les $\mathrm{Ext}$ à supports dans
$T$\footnote{La page écrit $\operatorname{Ext}_{T'}(X'; \ldots)$ et
$\operatorname{Ext}_T(F, \ldots)$ ; par excision, les $\mathrm{Ext}$ à supports
dans $T$, fermé dans $U$, se calculent sur $U$. Les indices de la formule
encadrée sont surchargés et en partie biffés sur la page.}. En effet, $p^{*}$
est exact et adjoint à gauche de $p_{*}$, qui transforme donc les injectifs en
injectifs ; $p$ est affine, de sorte que $p_{*}$ est exact sur les
quasi-cohérents, et $\Gamma_{T'} = \Gamma_T \circ p_{*}$. D'où
$\operatorname{Ext}^i_{T'}(U'; p^{*}F, p^{*}G) \simeq
\operatorname{Ext}^i_T(U; F, p_{*}p^{*}G)$ et
$p_{*}p^{*}G = \bigoplus_j G(j)$ ; enfin, $F$ étant cohérent et $U$
noethérien, $\operatorname{Ext}^i_T(U; F, -)$ commute aux sommes
directes\footnote{Les deux marginales de la page, « $G$ q.-coh. » et « $G$
q.-coh., $F$ de type fini », sont rattachées aux deux dernières égalités de la
chaîne ; la page passe par la résolution canonique $C(G')$ de Godement.}.
C'est la forme, avec $\mathrm{Ext}$ à supports, de ce qu'on appelle aujourd'hui
la correspondance de Serre-Grothendieck entre cohomologie locale graduée du
cône et cohomologie des tordus sur $\operatorname{Proj}$\footnote{Nom courant
dans les manuels (par exemple Brodmann et Sharp, \emph{Local cohomology}),
cité de mémoire ; le nom n'est pas sur la page.}.

La page 3 localise sur $Y$ et énonce le même isomorphisme pour les
$\underline{\operatorname{Ext}}$ relatifs à $Y$ sur le cône affine $C$, à
supports dans l'adhérence $\overline{T'}$ de $T'$, puis compare les
$\mathrm{Ext}$ sur $C$ à ceux du complété de $C$ le long de $\overline{T'}$,
par une flèche
\[
  \operatorname{Ext}^{\bullet}_{\widehat{T'}}(\widehat{C}; \widehat{F'}, \widehat{G'})
  \longrightarrow \operatorname{Ext}^{\bullet}_{\overline{T'}}(C; \overline{F'}, \overline{G'}),
\]
une bijection « par localisation » sur les
$H^p(\overline{Z'} - \overline{R'}, -)$, et un lemme des cinq qui devait
conclure. Presque tous les mots de liaison sont perdus, une marginale demande
de « justifier » une compatibilité, et la page s'arrête sur l'annonce d'un
lemme à dégager\footnote{$\overline{R'}$ n'est défini sur aucune page lisible ;
les croquis de la page 3 le placent comme le point où $\overline{Z'}$ rencontre
$C - \overline{U'}$. Nous ne reconstruisons pas l'argument.}. La page 4 met
face à face, terme à terme, deux paires de suites spectrales du local au
global :
\[
  H^p\bigl(U, \underline{\operatorname{Ext}}^q_{T}(F, G(j))\bigr) \Rightarrow
  \operatorname{Ext}^{p+q}_{T}(U; F, G(j)),
  \qquad
  H^p_T\bigl(U, \underline{\operatorname{Ext}}^q(F, G)(j)\bigr) \Rightarrow
  \operatorname{Ext}^{p+q}_{T}(U; F, G(j)),
\]
et leurs analogues sur le cône, avec une note sur leur compatibilité aux
graduations, et s'arrête sur « Cas particuliers lorsque $F =
\mathcal{O}_X$… ».

\pagerange{6}{15}
\section*{II. Dualité locale pour Modules non quasi cohérents (pages 6 à 15)}

\pagerange{6}{7}
\subsection*{Les $\mathrm{Ext}$ sur un schéma formel affine (pages 6 et 7)}

\textbf{Proposition 1.} Soient $A$ un anneau adique noethérien (séparé et
complet pour la topologie $J$-adique), $\mathfrak{X} = \operatorname{Spf} A$,
$M$, $N$ deux $A$-modules de type fini et $\mathcal{M} = M^{\Delta}$,
$\mathcal{N} = N^{\Delta}$ les $\mathcal{O}_{\mathfrak{X}}$-Modules cohérents
associés. Alors, les $\mathrm{Ext}$ étant pris dans la catégorie de tous les
$\mathcal{O}_{\mathfrak{X}}$-Modules,
\[
  \operatorname{Ext}^i(\mathfrak{X}; \mathcal{M}, \mathcal{N}) \simeq
  \operatorname{Ext}^i_A(M, N).
\]

\textbf{Corollaire 2.} $\underline{\operatorname{Ext}}^i(\mathcal{M},
\mathcal{N}) \simeq \operatorname{Ext}^i_A(M, N)^{\Delta}$.

La page donne deux démonstrations\footnote{Le début de l'énoncé de la
proposition 1 est d'une lecture très incertaine ; « adique », « séparé »,
« complet » sont des lectures douteuses que l'énoncé demande. La première
démonstration est entre crochets.}. La première passe par le corollaire : les
$\underline{\operatorname{Ext}}^i(\mathcal{M}, \mathcal{N})$ sont cohérents, la
cohomologie d'un Module cohérent sur $\operatorname{Spf} A$ est nulle en degré
$> 0$, donc la suite spectrale du local au global dégénère et
$\operatorname{Ext}^i(\mathfrak{X}; \mathcal{M}, \mathcal{N}) \simeq
H^0(\mathfrak{X}, \underline{\operatorname{Ext}}^i(\mathcal{M}, \mathcal{N}))$.
La seconde part d'une résolution de $M$ par des modules libres de type fini,
ou, autre procédé, de la platitude de $\mathfrak{X} \to \operatorname{Spec} A$ :
le complété commute à la formation des $\underline{\operatorname{Ext}}$ de
Modules cohérents, d'où
$\underline{\operatorname{Ext}}^i_{\mathcal{O}_{\mathfrak{X}}}(\mathcal{M},
\mathcal{N}) \simeq \underline{\operatorname{Ext}}^i(\widetilde{M},
\widetilde{N})^{\wedge} = \operatorname{Ext}^i_A(M, N)^{\Delta}$, c'est le
corollaire, et la proposition en prenant les sections.

La page 7 pose, pour $\mathfrak{X}$ préschéma formel localement noethérien,
$Y$ fermé et $\mathcal{F}$, $\mathcal{G}$ cohérents, la suite spectrale
\[
  H^p_Y\bigl(\mathfrak{X}, \underline{\operatorname{Ext}}^q(\mathcal{F},
  \mathcal{G})\bigr) \Rightarrow \operatorname{Ext}^{p+q}_Y(\mathfrak{X};
  \mathcal{F}, \mathcal{G}),
\]
et la compare à $\varprojlim_r \operatorname{Ext}^i_Y(\mathfrak{X};
\mathcal{F}, \mathcal{G}_r)$, où $\mathcal{G}_r$ est, on le suppose, la
réduction de $\mathcal{G}$ modulo la puissance $r$-ième d'un Idéal de
définition\footnote{La page note $\mathcal{G}_n$ sans le définir. La phrase
s'arrête sur « Notons d'ailleurs que $\operatorname{Ext}^{*}_Y(\mathfrak{X},
\mathcal{F}, \mathcal{G}_n)$ » et le reste du feuillet est blanc.}.

\pagerange{8}{8}
\subsection*{Notes éparses (page 8)}

Écrites à l'encre noire au verso d'un tapuscrit, elles touchent deux sujets.
D'abord trois isomorphismes entre $\mathrm{Ext}$ de schémas en groupes, à
valeurs dans ${}_{\infty}\mu$ d'un côté et dans $\mathbf{Q}/\mathbf{Z}$ de
l'autre, pour trois parties $U$, $Y \subset S$ ; leurs indices se croisent sur
la page, et nous ne les interprétons pas\footnote{La page écrit
$\underline{E}_{U/k}(G, {}_{\infty}\mu) \simeq \underline{E}_k(N^i_{U/k}\,G,
\mathbf{Q}/\mathbf{Z})$ et deux lignes analogues où les indices $Y/k$ et
$S/k$ sont croisés ; $N^i$ n'est pas défini.}.

Ensuite, pour un anneau de valuation discrète $V$, de corps résiduel $k$ et de
corps des fractions $K$, l'esquisse de la dualité locale en dimension $1$ :
la suite exacte
\[
  0 \to H^0(M) \to M \to M \otimes_V K \to H^1(M) \to 0
\]
est mise en regard, terme à terme, d'une suite d'$\mathrm{Ext}$ de $M$ dans
$V$, par des accouplements à valeurs dans $I$ --- en particulier
$M \otimes K$ contre $\operatorname{Hom}_K(M_K, K)$, et $H^1(M)$ contre
$\operatorname{Hom}(M, V)$ ---, avec en note « $\operatorname{Hom}(K, I)
\simeq K$ » et la formule
$\operatorname{Ext}_Y(M, A) \simeq \operatorname{Hom}(H^i_Y(X, M), I)$, dont
les degrés manquent\footnote{Les exposants des deux premiers $\mathrm{Ext}$ du
tableau ne se lisent pas ; la page note la base $A$ ici, $V$ juste au-dessus.
C'est le cas $n = 1$ des pages 10 et 11, pour des modules.}.

\pagerange{10}{10}
\subsection*{L'énoncé, et la réduction (page 10)}

Soient $A$ un anneau local noethérien \emph{régulier complet} de dimension $n$,
$X = \operatorname{Spec} A$, $x$ son point fermé, $I = H^n_x(\mathcal{O}_X)$,
et $F$ un $\mathcal{O}_X$-Module \emph{quelconque}, pas nécessairement
quasi-cohérent. Le cup-produit
\[
  H^i_x(F) \times \operatorname{Ext}^{n-i}(X; F, \mathcal{O}_X) \to
  H^n_x(\mathcal{O}_X) = I
\]
donne des homomorphismes
\[
  \varphi^{i}_F : \operatorname{Ext}^{n-i}(X; F, \mathcal{O}_X) \longrightarrow
  \operatorname{Hom}(H^i_x(F), I).
\]
La page affirme que ce sont des isomorphismes pour tout $i$ et tout $F$. Ce
n'est vrai que pour $n \leqslant 1$ ; les pages 14 et 15 le découvrent pour
$n = 2$\footnote{« pas nécessairement quasi-cohérent » est doublement souligné
sur la page ; la fin de la phrase « Je dis que ce sont des isomorphismes pour
tout $i$ » est en partie illisible. L'énoncé qui suit la réduction est celui de
la page ; ce qui est faux, c'est l'affirmation initiale, et nous le disons
ici plutôt que de la reproduire.}.

Tout $\mathcal{O}_X$-Module est quotient d'une somme de Modules
$\mathcal{O}_{U,X} = i_{!}(\mathcal{O}_U)$, où $i : U \to X$ est l'inclusion
d'un ouvert, qu'on peut prendre de la forme $X_f$ ; les deux membres sont des
foncteurs cohomologiques contravariants qui changent les sommes en produits.
La page ramène la question, par dévissage, aux $F = \mathcal{O}_{U,X}$, $U = X_f$,
pour lesquels $\operatorname{Ext}^{j}(X; \mathcal{O}_{U,X}, \mathcal{O}_X) =
H^j(U, \mathcal{O}_U)$ est nul pour $j > 0$ ($U$ est affine) et vaut $A_f$ pour
$j = 0$. L'énoncé équivaut donc, pour ces $U$, à
\[
  H^i_x(\mathcal{O}_{U,X}) = 0 \quad (i \neq n), \qquad
  D\bigl(H^n_x(\mathcal{O}_{U,X})\bigr) \simeq A_f
\]
par l'homomorphisme canonique. Pour $U = X$ c'est la dualité locale ordinaire
pour $A$ ; il reste les $U \neq X$.

\pagerange{10}{11}
\subsection*{Dimensions $0$ et $1$ : la dualité tient (pages 10 et 11)}

Pour $n = 0$, l'énoncé est une tautologie. Pour $n = 1$, $A$ est un anneau de
valuation discrète complet de corps des fractions $K$, $I = K/A$, et l'on a
l'autodualité $\operatorname{Hom}_A(K, K/A) \simeq K$. L'espace $X$ a deux
points, $x$ et le point générique, et la catégorie des $\mathcal{O}_X$-Modules
est équivalente à celle des triplets $(M_0, M_1, u)$ formés d'un $A$-module
$M_0 = F_x$, d'un $K$-espace vectoriel $M_1$ (la fibre générique) et d'une
application $A$-linéaire $u : M_0 \to M_1$. On a
\[
  H^0_x(F) = \operatorname{Ker} u, \qquad H^1_x(F) = \operatorname{Coker} u,
  \qquad H^0(X, F) = M_0, \quad H^i(X, F) = 0 \ (i > 0).
\]
La résolution injective $\mathcal{O}_X \to \underline{K}_X \to
\underline{K}_X/\mathcal{O}_X$ donne
\[
  \operatorname{Ext}^{\bullet}(X; F, \mathcal{O}_X) = H^{\bullet}\bigl(
  \operatorname{Hom}_K(M_1, K) \longrightarrow \operatorname{Hom}_A(M_0, K/A)
  \bigr),
\]
le complexe étant placé en degrés $0$ et $1$, la flèche étant $g \mapsto g
\circ u \bmod A$. Les homomorphismes de dualité deviennent
\[
  \varphi^0 : \operatorname{Ext}^0 \to \operatorname{Hom}(\operatorname{Coker} u,
  K/A), \qquad
  \varphi^1 : \operatorname{Ext}^1 \to \operatorname{Hom}(\operatorname{Ker} u,
  K/A),
\]
et ce sont des isomorphismes par les trois faits suivants :
\begin{enumerate}
\item[a)] tout homomorphisme d'un $K$-espace vectoriel dans $A$ est nul, son
image étant un sous-module divisible de $A$ ; donc $\varphi^0$ est injectif ;
\item[b)] $\operatorname{Ext}^1_A(K, A) = 0$, parce que $A$ est complet : dans
la suite exacte $\operatorname{Hom}(K, K) \to \operatorname{Hom}(K, K/A) \to
\operatorname{Ext}^1(K, A) \to 0$, la première flèche est surjective, car
$\operatorname{Hom}(K, K/A) \simeq \widehat{K} = K$. Tout homomorphisme
$M_1 \to K/A$ se relève donc en $M_1 \to K$ (automatiquement $K$-linéaire),
et $\varphi^0$ est surjectif ;
\item[c)] $K/A$ est injectif : tout homomorphisme $\operatorname{Ker} u \to K/A$
s'étend à $M_0$, et $\varphi^1$ est surjectif ; un homomorphisme $M_0 \to K/A$
nul sur $\operatorname{Ker} u$ se factorise par $u(M_0) \subset M_1$, s'étend
à $M_1$, se relève à $K$ par b), et vient donc de $\operatorname{Hom}(M_1, K)$
: $\varphi^1$ est injectif\footnote{La fin de c) ne se lit qu'en partie ; nous
avons rétabli l'argument dans l'ordre logique. L'insertion
« $\widehat{K} = K$ » et la remarque « $A$ complet » sont sur la page, au-dessus
de b).}.
\end{enumerate}

La dualité locale vaut donc, en dimension $1$, pour \emph{tous} les
$\mathcal{O}_X$-Modules --- à condition que $A$ soit complet, c'est le fait
b)\footnote{Pour un anneau de valuation discrète $A$, la nullité de
$\operatorname{Ext}^1_A(K, A)$ équivaut à la complétude de $A$ ; cité de
mémoire, sans prétendre à davantage en dimension supérieure. La page
15 pose la question de la nullité des $\operatorname{Ext}^i(K, A)$ pour $A$
local noethérien intègre complet.}.

\pagerange{12}{12}
\subsection*{Brouillon (page 12)}

La page 12, écrite au verso d'un tapuscrit, est un brouillon : des
$\operatorname{Ext}^i(K, A)$, et, en dimension $2$, la résolution injective
minimale de $A$ régulier,
\[
  0 \to A \to K \to \bigoplus_{\mathrm{ht}\,\mathfrak{p} = 1} K/A_{\mathfrak{p}}
  \to I \to 0,
\]
puis sa duale, qu'il écrit $0 \to K \to {\prod}^{\wedge}\, \widehat{K}_{\mathfrak{p}}
\to \operatorname{Hom}(K, I) \to 0$, en nommant le terme du milieu « produit
restreint », « i.e.\ adèles », avec, en regard, le carré
$A \to \prod \widehat{A}_{\mathfrak{p}}$, $K \to \prod
\widehat{K}_{\mathfrak{p}}$ et les mots « classes d'idèles », « classes
d'idéaux »\footnote{L'identification de la première suite comme résolution
injective minimale (le complexe de Cousin) est de cette lecture ; la page
l'appelle « résolution de $A$ ». Les indices sous les produits ne se lisent
pas. Nous ne vérifions pas l'exactitude de la seconde suite, que la page
n'affirme que par sa disposition.}.

\pagerange{14}{15}
\subsection*{Dimension $2$ : la dualité algébrique échoue (pages 14 et 15)}

Soient $n = 2$, $f \in \mathfrak{m}$ non nul, $U = X_f$, $F =
\mathcal{O}_{U,X}$, $X' = X - \{x\}$, $Y = V(f)$ et $Y' = Y \cap X'$, ensemble
fini et discret des points de hauteur $1$ contenant $f$, d'anneaux locaux
$\mathcal{O}_{x_i}$. Comme $H^0(X, F) = H^0(X', F) = 0$, on a $H^0_x(F) =
H^1_x(F) = 0$ ; $H^i_x(F) = 0$ pour $i > 2$ ; et, puisque $H^i(X, F) = 0$ pour
$i > 0$ ($X$ est le seul ouvert contenant $x$, de sorte que $\Gamma(X, -)$ est la fibre en $x$, foncteur exact), $H^2_x(F) \simeq H^1(X', F)$. La suite exacte $0 \to F \to
\mathcal{O}_{X'} \to G \to 0$, où $G$ est la restriction de $\mathcal{O}$ à
$Y'$, donne
\[
  0 \to A \xrightarrow{\ \alpha^0\ } \prod_i \mathcal{O}_{x_i}
  \xrightarrow{\ \partial\ } H^2_x(F) \xrightarrow{\ \alpha^1\ } I \to 0,
\]
avec $\prod_i \mathcal{O}_{x_i} = \varinjlim_g A[g^{-1}]$, $g$ parcourant les
éléments tels que $(f, g)$ soit $\mathfrak{m}$-primaire. En dualisant, on
obtient une suite exacte $0 \to A \to D(H^2_x(F)) \to D(\prod_i
\mathcal{O}_{x_i}) \to I \to 0$ et l'homomorphisme de dualité
$\varphi : A_f \to D(H^2_x(F))$. Le calcul explicite, $h \mapsto
[h][1/g]$, montre que le noyau de la flèche composée $A_f \to D(\prod_i
\mathcal{O}_{x_i})$ est l'image de $A$, et $\varphi$ est \emph{injectif} ; mais
son image est trop petite, et $\varphi$ \emph{n'est pas surjectif}. La raison,
dit la page, est qu'il faut prendre les \emph{duaux topologiques}, et non
seulement algébriques : mettre sur $\prod_i \mathcal{O}_{x_i}$ la topologie
produit des topologies naturelles des anneaux locaux, sur $H^2_x(F)$ la
topologie quotient, et ne retenir que les homomorphismes continus dans $I$
discret\footnote{La non-surjectivité est affirmée sur la page, avec la raison
qualitative qu'on vient de dire ; la page n'en donne pas d'autre preuve, et
nous n'en ajoutons pas. Le feuillet 15 est d'une écriture très rapide ; une
incise cerclée et un « N.B. » sur la complétion ne se lisent pas.}. La page
s'arrête sur « Question : dualité… » et la question des
$\operatorname{Ext}^i(K, A)$.

L'énoncé correct qui ressort des pages 10 à 15 est donc : la dualité locale
pour des $\mathcal{O}_X$-Modules arbitraires, avec des duaux algébriques, vaut
en dimension $\leqslant 1$ sur un anneau régulier complet et échoue en
dimension $2$ ; la version topologique est annoncée, non établie\footnote{Nous
ne connaissons pas d'énoncé publié qui porte ce titre, et n'en citons
aucun.}.

\pagerange{17}{34}
\section*{III. Notes antiques 1958 (pages 17 à 34)}

La chemise, page 16, porte « Notes antiques 1958 (Caractérisation des Coh.\
Mac., des Gorenstein etc) ». La date est la sienne, et le mot « antiques »
laisse penser que la chemise a été écrite plus tard que les notes ;
Montpellier date le dossier « à partir de 1967-1968 ». Rien sur les feuillets
ne permet de trancher\footnote{Si la date de la chemise est celle des notes,
elles précèdent le séminaire de Grothendieck sur la cohomologie locale (Harvard,
1961, publié par Hartshorne en 1967), l'article de Bass, \emph{On the
ubiquity of Gorenstein rings} (1963), et la théorie du module canonique de
Herzog et Kunz (1971). Les feuillets ne citent aucun de ces textes, et nous
n'en tirons aucune conclusion de priorité. Cité de mémoire.}.

\pagerange{17}{17}
\subsection*{Le plan (page 17)}

\textbf{§ 1. Les foncteurs $H^i_{\mathfrak{C}}$} --- foncteurs $A$-linéaires
(de modules et de modules topologiques) ; définition et premières propriétés
des $H^i_{\mathfrak{C}}(M)$ ; fonctorialité, y compris le cup-produit ; le cas
d'un anneau local, avec la caractérisation des $H^i$ et leurs relations avec
la profondeur ; rappels sur les modules de Cohen-Macaulay ; résultats
auxiliaires sur certains $\operatorname{Ext}^i_A(N, M)$, $M$ de
Cohen-Macaulay ; cas particuliers.

\textbf{§ 2. Foncteurs et modules dualisants (théorie locale). Le théorème de
dualité locale} --- définition et premières propriétés ; extension du foncteur
dualisant ; existence et unicité, exemples ; structure des $H^i(M)$ ; module
fondamental et théorème de dualité locale ; anneaux de Gorenstein.

\textbf{§ 3. Les foncteurs $H_i$ et les complexes dualisants.}

$H^i_{\mathfrak{C}}$ est la cohomologie locale à supports dans une famille
$\mathfrak{C}$\footnote{La lettre ronde $\mathfrak{C}$ n'est pas définie sur la
page ; c'est vraisemblablement une famille de supports, ou un idéal. Les
« foncteurs $H_i$ » et les « complexes dualisants » du § 3 n'apparaissent nulle
part ailleurs dans le dossier.}.

\pagerange{18}{18}
\subsection*{Le module fondamental (page 18)}

\textbf{Définition.} Soit $A$ un anneau local noethérien de dimension $n$. Un
$A$-module $\Omega$ est un \emph{module fondamental} s'il est de type fini et
si $D(\Omega) \simeq H^n(A)$.

Il est unique à isomorphisme près : $\widehat{\Omega} \simeq D(D(\Omega))
\simeq D(H^n(A))$ ne dépend que de $A$, et deux modules de type fini de
complétés isomorphes sont isomorphes. On le note $\Omega_A$. Il existe si $A$
est complet : on prend $D(H^n(A))$, de type fini par la dualité de Matlis
puisque $H^n(A)$ est artinien. Si $A$ est régulier, $\Omega_A = A$. Une
addition encadrée ajoute : si $A$ est de Cohen-Macaulay, module fondamental et
module dualisant coïncident\footnote{La notion de « module dualisant » d'un
anneau, distincte de $I$, devait être celle de l'article 1 du § 2, qui n'est
pas dans le dossier ; dans l'énoncé du théorème 4, « dualisant » est biffé et
remplacé par « fondamental ». Le début du feuillet est barré, et un bloc de
lignes biffées sépare la définition du théorème.}.

\textbf{Théorème 4.} Soit $A$ local noethérien de dimension $n$, muni d'un
module fondamental $\Omega$, et soit $H^n(\Omega) \simeq H^n(A) \otimes \Omega
\to I$ l'homomorphisme déduit d'un isomorphisme $H^n(A) \simeq D(\Omega)$. La
suite de l'énoncé ne se lit pas\footnote{Le bas du feuillet est d'une lecture
de plus en plus incertaine. D'après les corollaires de la page 20, l'énoncé
portait sur les accouplements $H^p(M) \times \operatorname{Ext}^{n-p}(M,
\Omega) \to I$ décrits ci-dessous.}.

\pagerange{19}{20}
\subsection*{La dualité locale (pages 19 et 20)}

Le foncteur $M \mapsto H^n(M)$ est exact à droite et commute aux limites
inductives, donc
\[
  H^n(M) \simeq H^n(A) \otimes_A M
\]
pour tout $M$ ; c'est le « complément au théorème 3 »\footnote{La page 19 est
barrée de deux longs traits, mais se lit en grande partie ; le théorème 3,
qu'elle complète, n'est pas dans le dossier. Un « Corollaire 2 » y est biffé
et repris en interligne par des fragments qui ne s'accordent pas (si
$\operatorname{Ext}^i(k, M) \neq 0$, alors $\operatorname{Ext}^j(k, M) = 0$
pour $j > i$ …) ; nous ne le restituons pas.}. Le cup-produit et
l'homomorphisme $H^n(\Omega) \to I$ donnent, pour tout $M$ de type fini, des
accouplements
\[
  H^p(M) \times \operatorname{Ext}^{n-p}(M, \Omega) \to H^n(\Omega) \to I,
  \qquad\text{d'où}\qquad
  \theta^p_M : H^p(M) \longrightarrow D\bigl(\operatorname{Ext}^{n-p}(M, \Omega)\bigr).
\]
En degré $n$, $\theta^n_M$ est la composée $H^n(M) \simeq H^n(A) \otimes M
\simeq D(\Omega) \otimes M \simeq D(\operatorname{Hom}(M, \Omega))$, donc un
isomorphisme. Les deux membres sont des foncteurs cohomologiques covariants de
$M$ ; par récurrence descendante sur $p$ et le lemme des cinq, appliqué à une
présentation $0 \to K \to L \to M \to 0$ avec $L$ libre, on obtient :

\emph{Pour un entier $k \leqslant n$, les $\theta^p_M$ sont des isomorphismes
pour tout $p \geqslant k$ et tout $M$ de type fini si et seulement si $H^i(A) =
0$ pour $k \leqslant i < n$.}

\textbf{Corollaire 1.} Les $\theta^p_M$ sont des isomorphismes pour tout $p$ et
tout $M$ de type fini si et seulement si $A$ est de Cohen-Macaulay (de
dimension $n$). C'est le \emph{théorème de dualité locale}\footnote{La page dit
« que les composantes de $A$ soient des anneaux de Coh.\ Mac.\ de dimension
$n$ », lecture très incertaine ; « composantes » est ajouté au-dessus de la
ligne. Pour $A$ local, la condition est que $A$ soit de Cohen-Macaulay : le
cas $M = A$ donne $H^p(A) \simeq D(\operatorname{Ext}^{n-p}(A, \Omega)) = 0$
pour $p < n$.}.

\textbf{Corollaire 2.} Si $A$ est de Cohen-Macaulay, $\operatorname{Ext}^i(k,
\Omega) = 0$ pour $i \neq n$ et $\operatorname{Ext}^n(k, \Omega) \simeq k$. En
effet, $\operatorname{Ext}^i(k, \Omega)$ est de longueur finie et son dual est
$H^{n-i}(k)$, nul si $i \neq n$ et égal à $k$ si $i = n$.

\textbf{Corollaire 3.} Si $A$ est de Cohen-Macaulay, $\Omega$ est un module de
Cohen-Macaulay de dimension $n$, $H^i(\Omega) = 0$ pour $i \neq n$,
$H^n(\Omega) \simeq I$, et
\[
  \operatorname{Ext}^i(\Omega, \Omega) = 0 \quad (i \neq 0), \qquad
  \operatorname{Hom}(\Omega, \Omega) \simeq A.
\]

\pagerange{21}{21}
\subsection*{Démonstration du corollaire 3, et deux corollaires de plus (page 21)}

Par le corollaire 2, $\operatorname{Ext}^i(k, \Omega) = 0$ pour $i < n$, donc
$\operatorname{prof} \Omega \geqslant n$, et $\Omega$ est de Cohen-Macaulay de
dimension $n$ ; $H^i(\Omega) = 0$ pour $i < n$. La dualité pour $M = \Omega$
donne $H^i(\Omega) \simeq D(\operatorname{Ext}^{n-i}(\Omega, \Omega))$, d'où
$\operatorname{Ext}^{j}(\Omega, \Omega) = 0$ pour $j > 0$, et, pour $i = n$,
$D(\operatorname{Hom}(\Omega, \Omega)) \simeq H^n(\Omega)$. Que $H^n(\Omega)
\simeq I = D(A)$ résulte du théorème 1 et du corollaire 1\footnote{Le théorème 1
n'est pas dans le dossier ; la phrase « D'après le Th.\ 1, le cor.\ 1 implique
que $H^n(\Omega)$ est dualisant » est encadrée sur la page.} ; d'où
$\operatorname{Hom}(\Omega, \Omega) \simeq A$, et l'argument de Nakayama de la
page 24 (ci-dessous) montre que l'isomorphisme est l'application canonique : tout
endomorphisme de $\Omega$ est la multiplication par un scalaire, déterminé de
façon unique ; a fortiori, $\Omega$ est fidèle.

\textbf{Corollaire 4.} Si $A$ est de Cohen-Macaulay, $\Omega_A$ est de dimension
injective $n$ : $\operatorname{Ext}^i(M, \Omega) = 0$ pour $i > n$ et tout $M$.
Il suffit de le voir pour $M$ de type fini, et c'est la dualité :
$\operatorname{Ext}^i(M, \Omega)$ est dual de $H^{n-i}(M) = 0$.

\textbf{Corollaire 5.} Soient $A$ de Cohen-Macaulay de dimension $n$ et $M$ de
type fini. Alors
\[
  \operatorname{prof} M > k \iff \operatorname{Ext}^i_A(M, \Omega) = 0 \ \text{pour}\
  i \geqslant n - k,
\]
et, si $\dim M = d$, $M$ est de Cohen-Macaulay si et seulement si
$\operatorname{Ext}^i_A(M, \Omega) = 0$ pour $i \neq n - d$\footnote{Le
corollaire 5 est écrit le long du bord gauche, feuille tournée. La page écrit
« $\mathrm{codim}\, M \leq k \Leftrightarrow \mathrm{Ext}^i_A(M, \Omega) = 0$
pour $i \geq n - k$ » (le signe $\geq$ est incertain) et « $\mathrm{Ext}^i = 0$
pour $i \neq k$ ». Par la dualité, $\operatorname{Ext}^{n-p}(M, \Omega)$ est dual
de $H^p(M)$, non nul pour $p = \operatorname{prof} M$ et $p = \dim M$, nul pour
$p < \operatorname{prof} M$ et $p > \dim M$ : l'inégalité sur la profondeur est
dans l'autre sens, et l'exposant est $n - d$.}.

\pagerange{22}{25}
\subsection*{L'accouplement $\mathfrak{D}$ et l'annulateur de $\Omega$ (pages 22 à 25)}

On ne suppose plus $A$ de Cohen-Macaulay. Soit $\Omega$ un module fondamental
pour $A$, muni d'un isomorphisme $\mathfrak{D}_1 : H^n(A) \xrightarrow{\sim}
D(\Omega)$, c'est-à-dire d'un accouplement
\[
  \mathfrak{D} : H^n(\Omega) = H^n(A) \otimes \Omega \longrightarrow I,
  \qquad \mathfrak{D} \in D(H^n(\Omega)),
\]
ou encore d'un isomorphisme $\mathfrak{D}_2 : \widehat{\Omega} \xrightarrow{\sim}
D(H^n(A))$ (ici le chapeau est le complété). Pour $M$ de type fini on en tire
\[
  \mathfrak{D}_{1,M} : H^n(M) \xrightarrow{\sim} D(\operatorname{Hom}(M, \Omega)),
  \qquad
  \mathfrak{D}_{2,M} : \operatorname{Hom}(M, \Omega)^{\wedge} \xrightarrow{\sim}
  D(H^n(M)),
\]
qui se réduisent à $\mathfrak{D}_1$, $\mathfrak{D}_2$ pour $M = A$. Pour
$M = \Omega$, l'identité de $\operatorname{Hom}(\Omega, \Omega)^{\wedge}$
correspond à $\mathfrak{D} \in D(H^n(\Omega))$ ; la transposée de $\mathfrak{D}$
est donc l'homomorphisme $\widehat{A} \to \operatorname{Hom}(\Omega,
\Omega)^{\wedge}$ qui envoie $\lambda$ sur l'homothétie de rapport
$\lambda$.

\textbf{Proposition 9.} Supposons $A$ complet\footnote{L'hypothèse est
ajoutée : la page renvoie au théorème 3, absent, pour l'annulateur de
$H^n(A)$. Pour $A$ complet, ou plus généralement quotient d'un anneau de
Gorenstein, l'énoncé est un résultat connu d'Aoyama (1980) ; sans hypothèse, il
vaut pour $\widehat{A}$. Cité de mémoire.}. L'annulateur de $\Omega$ est
l'intersection $\mathfrak{N}$ des composantes primaires de $0$ dans $A$ dont le
premier associé est de dimension $n$ ; $\Omega$ est aussi un module fondamental
pour $A/\mathfrak{N}$ ; les idéaux premiers associés à $\Omega$ sont les
idéaux premiers de dimension $n$ de $A$, et en ces premiers
$\mathrm{lg}_{\mathfrak{p}}(\Omega) = \mathrm{lg}_{\mathfrak{p}}(A)$\footnote{La
dernière égalité est en marge de l'énoncé, encadrée d'un trait vertical.}.

L'annulateur de $\Omega$ est celui de $D(\Omega) = H^n(A)$. La suite exacte
$0 \to \mathfrak{N} \to A \to A/\mathfrak{N} \to 0$, avec $\dim \mathfrak{N} <
n$, donc $H^n(\mathfrak{N}) = 0$, donne $H^n(A) \simeq H^n(A/\mathfrak{N})$,
d'où la deuxième assertion ; $\mathfrak{N}$ est le plus grand idéal de $A$ de
dimension $< n$.

\textbf{Corollaire 1.} L'image de $\mathfrak{D}$ est l'annulateur de
$\mathfrak{N}$ dans $I$, et $\operatorname{Coker} \mathfrak{D} \simeq
D(\mathfrak{N})$. Les conditions suivantes sont équivalentes : (i)
$\mathfrak{D}$ est surjectif ; (ii) $\mathfrak{N} = 0$, i.e.\ tous les idéaux
premiers associés à $0$ dans $A$ sont de dimension $n$ ; (iii) $\Omega$ est
fidèle. En effet $\mathfrak{N}$ est le noyau de $A \to \operatorname{Hom}(\Omega,
\Omega)$, transposée de $\mathfrak{D}$.

\textbf{Corollaire 2.} $\operatorname{Ker} \mathfrak{D} \simeq
D\bigl(\operatorname{Hom}(\Omega, \Omega)/A \cdot \mathrm{id}_{\Omega}\bigr)$,
et les conditions suivantes sont équivalentes : (i) $\mathfrak{D}$ est
injectif ; (i bis) $H^n(\Omega)$ est comonogène ; (ii) tout endomorphisme de
$\Omega$ est une homothétie ; (ii bis) $\operatorname{Hom}(\Omega, \Omega)$ est
monogène. La formule vient de la dualité ; (i) $\Leftrightarrow$ (ii bis)
$\Leftrightarrow$ (i bis), (ii) $\Rightarrow$ (ii bis) est trivial, et (ii bis)
$\Rightarrow$ (ii) parce que, si $\operatorname{Hom}(M, M)$ est monogène pour
$M$ de type fini, le lemme de Nakayama montre que $\mathrm{id}_M$ en est un
générateur.

\textbf{Corollaire 3.} Équivalent : (i) $\mathfrak{D}$ est bijectif ; (ii)
$H^n(\Omega) \simeq I$ ; (iii) $A \to \operatorname{Hom}(\Omega, \Omega)$ est
un isomorphisme ; (iv) $\operatorname{Hom}(\Omega, \Omega) \simeq A$. Alors
tous les premiers associés de $A$ sont de dimension $n$.

\textbf{Corollaire 4.} Les conditions du corollaire 3 sont remplies si (a) $A$
est de Cohen-Macaulay --- c'est le corollaire 3 du théorème 4, page 20 ---, ou
si (b) $A$ est normal. Dans le cas (b), $A$ est intègre, la condition du
corollaire 1 est remplie, et $\Omega$, sans torsion et de rang $1$, a son
anneau d'endomorphismes contenu dans le corps des fractions $K$ ; fini sur $A$,
il est contenu dans la clôture intégrale, c'est-à-dire dans $A$.

\textbf{Question.} Les conditions du corollaire 2 sont-elles toujours remplies,
au moins quand les premiers associés à $0$ sont tous de dimension $n$ ? La page
répond non, par un anneau local intègre de dimension $2$ dont les localisés de
hauteur $1$ sont unibranches, pour lequel
\[
  \operatorname{Hom}_A(\Omega, \Omega) = \bigcap_{\mathrm{ht}\,\mathfrak{p} = 1}
  A_{\mathfrak{p}} .
\]
Cette formule est juste pour un anneau local intègre de dimension $2$ : le
second membre est le plus petit anneau intermédiaire vérifiant la condition
$(S_2)$ de Serre\footnote{On l'appelle aujourd'hui la $S_2$-ification de $A$ ;
le nom est de nous. La page continue : « $=$ clôture intégrale de $A$ », écrit
au-dessus de « normalisé » biffé. Cette seconde égalité demande que les
$A_{\mathfrak{p}}$ de hauteur $1$ soient des anneaux de valuation discrète ;
« unibranche » ne suffit pas. Le contre-exemple tient dès que le second membre
diffère de $A$, c'est-à-dire dès que $A$ n'est pas de profondeur $2$.}. Elle
diffère de $A$ dès que $A$ est de profondeur $1$ ; par exemple pour $A =
k[[s^4, s^3t, st^3, t^4]]$, elle donne $k[[s^4, s^3t, s^2t^2, st^3, t^4]]$, qui
n'est pas monogène sur $A$ : la réponse « non » est juste\footnote{L'exemple
est de cette lecture ; la page n'en nomme aucun.}.

\textbf{Proposition 10.} Soit $x \in \mathfrak{m}$ non diviseur de zéro dans
$A$, tel que $x H^{n-1}(A) = H^{n-1}(A)$. Alors $\Omega_A/x\Omega_A$ est un
module fondamental pour $A/xA$\footnote{La page ne dit pas $x \in
\mathfrak{m}$, que l'énoncé demande.}.

\pagerange{26}{28}
\subsection*{Les anneaux quotients, et les modules $\Omega_A^{(i)}$ (pages 26 à 28)}

\emph{Démonstration de la proposition 10.} $\dim A/xA = n-1$, et la suite
exacte de cohomologie locale de $0 \to A \xrightarrow{x} A \to A/xA \to 0$
donne, compte tenu de l'hypothèse, $H^{n-1}(A/xA) \simeq \operatorname{Ker}(x
\mid H^n(A))$ ; comme $H^n(A) \simeq D(\Omega)$, ce noyau est $D(\Omega/x\Omega)$.

\textbf{Corollaire.} Si $A$ est de Cohen-Macaulay et $x_1, \ldots, x_p$ une
partie d'un système de paramètres, $\Omega_A/(x_1, \ldots, x_p)\Omega_A$ est
fondamental pour $A/(x_1, \ldots, x_p)$.

Supposons maintenant $A = A'/\mathfrak{a}$, quotient d'un anneau local de
Cohen-Macaulay $A'$ de dimension $n'$ muni d'un module fondamental
$\Omega_{A'}$. Un $A$-module est un $A'$-module, avec la même cohomologie
locale, et $I$ s'identifie au sous-module de $I'$ annulé par $\mathfrak{a}$. La
dualité locale sur $A'$ donne, pour tout $A$-module $M$ de type fini,
\[
  (1) \qquad H^i(M) \simeq D\bigl(\operatorname{Ext}^{n'-i}_{A'}(M, \Omega_{A'})\bigr),
\]
en particulier $H^n(A) \simeq D(\operatorname{Ext}^{n'-n}_{A'}(A,
\Omega_{A'}))$, d'où
\[
  (2\ \text{bis}) \qquad \Omega_A \simeq \operatorname{Ext}^{n'-n}_{A'}(A, \Omega_{A'}) .
\]
\footnote{L'exposant de (1) est noirci devant « $-i$ » ; (1 bis), qui portait
$\operatorname{Ext}^{n'-i}_{A'}$, est biffé. Dans (2), (2 bis) et (13), la page
écrit $\operatorname{Ext}_{A'}(M, \Omega_{A'})$ là où il faut $A$ : nous
corrigeons.}

Pour $\mathfrak{p}$ premier de $A$, image d'un premier $\mathfrak{p}' \supset
\mathfrak{a}$ de $A'$, la formation de (2 bis) commute à la localisation :
$(\Omega_A)_{\mathfrak{p}} \simeq \operatorname{Ext}^{n'-n}_{A'_{\mathfrak{p}'}}
(A_{\mathfrak{p}}, (\Omega_{A'})_{\mathfrak{p}'})$. Si
$(\Omega_{A'})_{\mathfrak{p}'}$ est fondamental pour $A'_{\mathfrak{p}'}$ --- par
exemple si $A'$ est régulier, car alors $\Omega_{A'} = A'$ et
$A'_{\mathfrak{p}'}$ est régulier (Serre) ---, la formule (2 bis) appliquée à
$A_{\mathfrak{p}} = A'_{\mathfrak{p}'}/\mathfrak{a}A'_{\mathfrak{p}'}$ montre
que $(\Omega_A)_{\mathfrak{p}}$ est fondamental pour $A_{\mathfrak{p}}$, pourvu
que l'exposant soit le bon, c'est-à-dire que $\dim A'_{\mathfrak{p}'} - \dim
A_{\mathfrak{p}} = n' - n$. C'est le cas quand $\mathfrak{p}$ contient un
premier minimal de dimension $n$, c'est-à-dire quand $\mathfrak{p} \in
\operatorname{Supp} \Omega_A$ (proposition 9)\footnote{La page conclut
« $(\Omega_A)_{\mathfrak{p}}$ est un module dualisant pour $A_{\mathfrak{p}}$ »,
où l'on attend « fondamental », et ne pose pas de condition sur
$\mathfrak{p}$. Elle écrit aussi « $\mathfrak{p} \simeq A'/\mathfrak{p}'$ » et
« $M_{\mathfrak{p}^*}$ », tels quels.}.

\textbf{Proposition 11.} Soit $A$ un quotient d'un anneau local régulier.
Alors $A$ a un module fondamental $\Omega$, et pour tout premier $\mathfrak{p}
\in \operatorname{Supp} \Omega$ --- en particulier pour tout $\mathfrak{p}$ si
$A$ est équidimensionnel ---, $\Omega_{\mathfrak{p}}$ est fondamental pour
$A_{\mathfrak{p}}$\footnote{La page dit « si $\mathfrak{p} \subset A$ est
premier » sans restriction ; c'est faux hors du support. Pour $A =
k[[x,y,z]]/(xy, xz)$, réunion d'un plan et d'une droite, $\Omega_A$ est porté
par le plan ; au point générique $\mathfrak{p} = (y, z)$ de la droite,
$(\Omega_A)_{\mathfrak{p}} = 0$, alors que $A_{\mathfrak{p}}$ est un corps,
module fondamental de lui-même. L'exemple est de cette lecture.}.

Une note demande une démonstration directe, sans passer par
l'interprétation\footnote{« Une démonstration directe … question de définition,
sans interprétation » : lecture partielle.}. Puis :

\textbf{Lemme.} Avec $A'$ régulier de dimension $n'$ et $M$ de type fini sur
$A$, on a $\operatorname{prof} M \geqslant k$ si et seulement si
$\operatorname{Ext}^{n'-i}_{A'}(M, \Omega_{A'}) = 0$ pour $0 \leqslant i <
k$\footnote{La page écrit « pour que $\mathrm{codim}\, M \leq k$ »
(« codim » est une lecture incertaine). Par (1), la condition dit
$H^i(M) = 0$ pour $i < k$, c'est-à-dire $\operatorname{prof} M \geqslant k$.}.

\textbf{Corollaire.} $A$ est de Cohen-Macaulay si et seulement si
$\operatorname{Ext}^{n'-n+i}_{A'}(A, \Omega_{A'}) = 0$ pour $i > 0$.

Posons
\[
  (13) \qquad \Omega_A^{(i)} = \operatorname{Ext}^{n'-n+i}_{A'}(A, \Omega_{A'}),
  \qquad\text{de sorte que}\qquad
  (13\ \text{bis}) \qquad D(\Omega_A^{(i)}) \simeq H^{n-i}(A)
\]
à isomorphisme non canonique près. Alors $\Omega_A^{(0)} = \Omega_A$, et
$\Omega_A^{(i)} = 0$ pour tout $i \geqslant 1$ si et seulement si $A$ est de
Cohen-Macaulay. Ce sont les \emph{modules de déficience} de
$A$\footnote{P.~Schenzel, \emph{Dualisierende Komplexe in der lokalen Algebra
und Buchsbaum-Ringe} (1982), qui note $K^{j}(A) =
\operatorname{Ext}^{n'-j}_{A'}(A, A')$ ; ici $\Omega_A^{(i)} = K^{n-i}(A)$. Cité
de mémoire. La page esquisse une « analogie » entre crochets, illisible. Les
exposants de (13) et (13 bis) sont repassés à l'encre noire.}.

\pagerange{29}{29}
\subsection*{Une suite spectrale (page 29)}

\textbf{Lemme 6.} Soient $A' \to A$ un homomorphisme d'anneaux, $N$ un
$A$-module et $M$ un $A'$-module. Il y a une suite spectrale
\[
  E_2^{p,q} = \operatorname{Ext}^p_A\bigl(N, \operatorname{Ext}^q_{A'}(A, M)\bigr)
  \Rightarrow \operatorname{Ext}^{p+q}_{A'}(N, M).
\]
C'est la suite spectrale du foncteur composé $\operatorname{Hom}_{A'}(N, -) =
\operatorname{Hom}_A(N, -) \circ \operatorname{Hom}_{A'}(A, -)$, le second
foncteur transformant injectifs en injectifs (il est adjoint à droite de la
restriction des scalaires, qui est exacte)\footnote{Plusieurs $A$ de la page
sont repassés sur une première lettre, qui semble un $B$ ; il reste
$\operatorname{Hom}_{A'}(B, M)$, qu'on lit $\operatorname{Hom}_{A'}(A, M)$.}.

\textbf{Proposition 12.} Soit $A$ un quotient d'un anneau local de
Cohen-Macaulay $A'$ muni d'un module fondamental. Pour tout $A$-module $M$ de
type fini, il y a une suite spectrale
\[
  E_2^{p,q} = \operatorname{Ext}^p_A(M, \Omega_A^{(q)}) \Rightarrow
  D\bigl(H^{n-p-q}(M)\bigr)
\]
(à complétion près), et la forme encadrée sur la page, $D(\operatorname{Ext}^p_A(M,
\Omega_A^{(q)})) \Rightarrow H^{n-p-q}(M)$, s'en déduit en dualisant. C'est le
lemme 6 avec $N = M$ et $\Omega_{A'}$ à la place de $M$, suivi de (1). Elle
« explique, en le généralisant, le théorème de dualité locale » : si $A$ est de
Cohen-Macaulay, seule la ligne $q = 0$ subsiste\footnote{L'hypothèse sur $A$
est écrite au-dessus de l'énoncé ; l'indice de $H$ dans la formule encadrée est
précédé d'un signe noirci, illisible, que nous rétablissons comme $n - \cdot$.
La dernière phrase, sur le cas Cohen-Macaulay, est de cette lecture.}.

En marge, le corollaire 1 : si $x \in \mathfrak{m}$ n'est diviseur de zéro ni
dans $A$ ni dans $M$, on a $\operatorname{Ext}^p_A(N, M) \simeq
\operatorname{Ext}^{p-1}_{A/xA}(N, M/xM)$ pour tout $A/xA$-module $N$ ---
c'est le lemme 6 appliqué à $A \to A/xA$, et c'est ce qu'on appelle aujourd'hui
le lemme de Rees\footnote{La page dit « pour tout $N$ », sans préciser que $N$
est un $A/xA$-module ; le nom est de nous. Un corollaire 2 et une annonce en
oblique au pied de la marge ne se lisent pas.}.

\pagerange{30}{34}
\subsection*{6. Modules fondamentaux des anneaux de Cohen-Macaulay ; anneaux de Gorenstein (pages 30 à 34)}

Dans ce paragraphe, $M$ est un $A$-module de type fini et $n = \dim A >
0$\footnote{La page dit « $M$ un $A$-module, $n$ un entier $> 0$ ». Pour
$n$ quelconque, l'implication c') $\Rightarrow$ a') ci-dessous, qui utilise
$H^{n+1}(M) = 0$, demande $n \geqslant \dim M$.}. La page suppose
$\operatorname{Ext}^{n-1}(k, M) = 0$ ; ses démonstrations utilisent
l'isomorphisme $\operatorname{Ext}^n(N, M) \simeq \operatorname{Hom}(N,
H^n(M))$ pour $N$ de longueur finie, qui demande $\operatorname{Ext}^i(k, M) =
0$ pour \emph{tout} $i < n$, c'est-à-dire $\operatorname{prof} M \geqslant n$ :
c'est l'hypothèse sous laquelle nous énonçons\footnote{Pour $M$ de type fini et
$\operatorname{Ext}^n(k, M) \neq 0$, la seule condition
$\operatorname{Ext}^{n-1}(k, M) = 0$ entraîne en fait $\operatorname{prof} M
\geqslant n$, mais par des résultats bien postérieurs : la non-nullité des
nombres de Bass $\mu_i(M) = \dim_k \operatorname{Ext}^i(k, M)$ pour
$\operatorname{prof} M \leqslant i \leqslant \operatorname{inj.dim} M$
(Fossum, Foxby, Griffith et Reiten, 1975 ; Roberts) et l'égalité
$\operatorname{inj.dim} M = \operatorname{prof} A$ quand elle est finie (Bass,
1963). Cité de mémoire. La remarque de la page 32 et le Problème C de la page
37 montrent que la question était ouverte pour lui.}. Pour $M \neq 0$, elle
signifie que $M$ est de Cohen-Macaulay de dimension $n$.

\textbf{Proposition 13.} Sous ces hypothèses, les conditions suivantes sont
équivalentes :
\begin{enumerate}
\item[a')] $\operatorname{Ext}^{n+1}(k, M) = 0$ ;
\item[b')] $N \mapsto \operatorname{Ext}^n(N, M)$ est exact sur les modules de
longueur finie ;
\item[c')] $H^n(M)$ est injectif.
\end{enumerate}
a') $\Rightarrow$ b') par dévissage ; b') $\Leftrightarrow$ c') parce que
$\operatorname{Ext}^n(N, M) \simeq \operatorname{Hom}(N, H^n(M))$ et que
$H^n(M)$ est de $\mathfrak{m}$-torsion ; c') $\Rightarrow$ a') : l'exactitude
de $\operatorname{Ext}^n(-, M)$ rend $\operatorname{Ext}^{n+1}(-, M)$ exact à
gauche, donc $\operatorname{Ext}^{n+1}(k, M) \to H^{n+1}(M) = 0$ est
injectif\footnote{Une première version, sur la même page barrée d'une croix,
portait sous le nom de théorème 5 puis de proposition 13 l'équivalence de
a) $\operatorname{Ext}^n(k, M) = k$ et $\operatorname{Ext}^{n+1}(k, M) = 0$,
b) $N \mapsto \operatorname{Ext}^n(N, M)$ dualisant, c) $H^n(M) \simeq I$, avec
une condition sur les longueurs ; elle revient dans le corollaire 2. La version
gardée, en pied de page, porte « Proposition 13 » sur un « Lemme 7 » biffé.}.

\textbf{Corollaire 1.} Sous ces conditions, $H^n(M) \simeq I^m$ pour un entier
$m$, $\operatorname{Ext}^n(k, M) \simeq k^m$, $\operatorname{Ext}^i(k, M) = 0$
pour $i > n$, et $\mathrm{lg}_{\mathfrak{p}}(M) = m\,
\mathrm{lg}_{\mathfrak{p}}(A)$ pour tout premier minimal $\mathfrak{p}$ de
dimension $n$ (de $\widehat{A}$ si $A$ n'est pas complet)\footnote{L'énoncé
est très repris ; plusieurs incises entre crochets ne se lisent pas, et le
« $\sqrt{A}$ » de la page est une lecture incertaine.}.

\textbf{Proposition 14.} Sous ces conditions, $A$ est de Cohen-Macaulay et
$\widehat{M} \simeq \Omega_{\widehat{A}}^m$. Une « Remarque » qui la précède, serrée et en grande partie illisible,
l'annonçait déjà pour $\Omega^m$, en notant que l'hypothèse sur
$\operatorname{Ext}^{n-1}$ serait redondante dans b'). La page procède par
récurrence sur $n$, triviale pour $n = 0$, et s'arrête\footnote{« Voir les \ldots » suivi de
tirets. L'énoncé est vrai : $M$ est de Cohen-Macaulay de dimension $n$ et de
dimension injective finie, égale à $n$ ; elle est égale à $\operatorname{prof}
A$ (Bass, 1963), donc $A$ est de Cohen-Macaulay, et un tel module est une
somme de copies du module canonique (voir Bruns et Herzog,
\emph{Cohen-Macaulay rings}, 3.3.28). Cité de mémoire. Le cas $m = 1$ est
démontré pages 38 et 39.}.

En marge, le \textbf{corollaire 2} : jointes à $0 <
\mathrm{lg}_{\mathfrak{p}}(M) \leqslant \mathrm{lg}_{\mathfrak{p}}(A)$, les
conditions précédentes équivalent à a) $\operatorname{Ext}^n(k, M) = k$ et
$\operatorname{Ext}^{n+1}(k, M) = 0$ ; b) $N \mapsto \operatorname{Ext}^n(N, M)$
est dualisant ; c) $H^n(M) \simeq I$ --- c'est le cas $m = 1$. Et le
\textbf{corollaire 3} : elles sont entraînées par d) $\operatorname{Ext}^n(k, M)
= k$ et $M$ fidèle, ou d bis) $\operatorname{Ext}^n(A/\mathfrak{q}, M)$
comonogène et $M$ fidèle.

Page 32, le corollaire 1 de la proposition 15 (numéro biffé) : pour $M$ de
Cohen-Macaulay de dimension $n$, les conditions a), b), c) du corollaire 2 et
d), d bis) du corollaire 3 sont équivalentes --- elles disent que $A$ est de
Cohen-Macaulay et $M \simeq \Omega_A$ ---, et équivalentes à

f) $M/(x_1M + \cdots + x_nM)$ est comonogène, pour un système de paramètres
$x_1, \ldots, x_n$ de $M$ ;

car $\operatorname{Ext}^n(k, M) \simeq \operatorname{Hom}(k, M/(x_1, \ldots,
x_n)M)$. Une remarque se demande si « $M$ de Cohen-Macaulay » peut y être
remplacé par $\operatorname{Ext}^{n-1}(k, M) = 0$.

\textbf{Proposition 15.} Les conditions suivantes sont équivalentes :
\begin{enumerate}
\item[(a)] $A$ est un module fondamental pour lui-même ;
\item[(b)] $H^n(A) \simeq I$ ;
\item[(c)] $H^n(A)$ est injectif ;
\item[(d)] $H^n(A)$ est comonogène, et tous les idéaux premiers associés de
$\widehat{A}$ sont de dimension $n$.
\end{enumerate}
(a) dit $D(A) = I \simeq H^n(A)$, c'est (b) ; (b) $\Rightarrow$ (c) ; (c)
donne $H^n(A) \simeq I^m$, donc $A^m$ fondamental, et la proposition 9 donne
$\mathrm{lg}_{\mathfrak{p}}(A^m) = \mathrm{lg}_{\mathfrak{p}}(A)$, d'où $m = 1$
et (a). (a) $\Rightarrow$ (d) par la proposition 9 ; réciproquement, sous (d),
$H^n(A)$ se plonge dans $I$, donc $\widehat{A}$ se surjecte sur
$D(H^n(A)) = \Omega_{\widehat{A}}$, avec pour noyau l'annulateur
$\mathfrak{N}$ de $\Omega_{\widehat{A}}$, nul par la proposition 9\footnote{La
page écrit en (d) « $A$ est équidimensionnel ». C'est trop faible : pour $A =
k[[x,y]]/(x^2, xy)$, de dimension $1$, qui n'a qu'un premier minimal, $(x)$,
on a $\mathfrak{N} = (x)$, $\Omega_A \simeq k[[y]]$, monogène, donc $H^1(A) =
D(\Omega_A)$ comonogène, et pourtant $A \not\simeq \Omega_A$. Il faut que
$\widehat{A}$ n'ait pas de premier associé immergé ni de composante de
dimension $< n$. L'exemple est de cette lecture.}.

On appelle aujourd'hui ces anneaux \emph{quasi-Gorenstein}\footnote{Le nom,
postérieur, est d'Aoyama et de Platte et Storch (vers 1977-1980) ; cité de
mémoire.}. La page 33 avait d'abord posé : « \emph{Définition 2.} Un tel $A$
est dit anneau de Gorenstein », puis « \emph{Th.\ 5.} Un anneau de Gorenstein
est Coh.\ Mac. », et commencé une récurrence par un non-diviseur de zéro $x$ ;
tout est biffé, définition, théorème et démonstration. Le théorème biffé est
faux : le cône affine (complété au sommet) sur une surface abélienne
projectivement normale est quasi-Gorenstein --- son module canonique est
l'anneau lui-même, parce que le fibré canonique de la surface est trivial ---
et n'est pas de Cohen-Macaulay, parce que $H^2_{\mathfrak{m}}$ de l'anneau
contient $H^1$ de la surface, qui est non nul\footnote{L'exemple est classique
et de cette lecture ; la page n'en donne aucun.}.

\textbf{Proposition 16.} Si $A$ est de Cohen-Macaulay, les conditions de la
proposition 15 équivalent aussi à : b') $N \mapsto \operatorname{Ext}^n(N, A)$
est dualisant ; c') il est exact ; c'') $\operatorname{Ext}^{n+1}(k, A) =
\operatorname{Ext}^{n-1}(k, A) = 0$ ; d) $\operatorname{Ext}^n(k, A) = k$. Une
marge affirme qu'il est facile de voir que b') et c') entraînent déjà que $A$
est de Cohen-Macaulay ; la page ne le montre pas, et nous ne l'avons pas
vérifié\footnote{La fin de la page 33 est surchargée : lettres reprises,
flèches, et une longue boucle qui déplace les conditions ; l'ordre est celui
de la page 34.}.

\textbf{Corollaire 1.} Elles équivalent encore à : e') $\operatorname{Ext}^n(A/
\mathfrak{q}, A)$ est comonogène ; e'') $A/(x_1, \ldots, x_n)A$ est comonogène,
c'est-à-dire autodual.

\textbf{Corollaire 2.} Si $A$ est quotient d'un anneau régulier $A'$ de
dimension $n'$, elles équivalent à $\operatorname{Ext}^{n'-n}_{A'}(A, A') \simeq
A$.

\textbf{Remarque.} « Il se peut que dans toutes les conditions, l'hypothèse
Cohen-Macaulay soit surabondante. » Elle ne l'est pas pour (a)--(d), on vient
de le voir ; elle l'est pour c''), par les résultats sur les nombres de Bass
déjà cités\footnote{Si $\operatorname{Ext}^{n+1}(k, A) = 0$ avec
$\operatorname{prof} A \leqslant n$, la dimension injective de $A$ est finie,
donc $A$ est de Gorenstein, en particulier de Cohen-Macaulay. Cité de
mémoire.}.

\textbf{Définition 2.} Un \emph{anneau de Gorenstein} est un anneau local de
Cohen-Macaulay qui est un module fondamental pour lui-même. Exemple : les
anneaux réguliers.

\textbf{Proposition 17.} (i) Un localisé d'un anneau de Gorenstein (supposé
quotient d'un anneau régulier) est de Gorenstein. (ii) De même pour le
quotient par l'idéal engendré par une partie d'un système de paramètres.

\textbf{Corollaire} (inachevé sur la page) : le quotient d'un anneau régulier
par un idéal engendré par une partie d'un système de paramètres --- une
intersection complète --- est de Gorenstein\footnote{La phrase s'arrête sans
verbe ; « partie d'un » est ajouté au-dessus de la ligne. La conclusion est
celle que (ii) donne.}.

\pagerange{35}{39}
\section*{IV. Récapitulation, théorème, problèmes (pages 35 à 39)}

\pagerange{35}{36}
\subsection*{Les conditions (pages 35 et 36)}

D'une écriture plus posée, la page 35 fixe $M$ de type fini sur $A$, un idéal
de définition $\mathfrak{q}$, un système de paramètres $x_1, \ldots, x_n$, un
premier $\mathfrak{p}$ de dimension $n$, $N$ variable de longueur finie, et
dresse la liste :
\begin{enumerate}
\item[(i)] $\operatorname{Ext}^{n-1}(k, M) = 0$ et l'une des conditions a)
$\operatorname{Ext}^{n+1}(k, M) = 0$ et $\operatorname{Ext}^n(k, M) \simeq k$ ;
b) $\operatorname{Ext}^n(-, M)$ dualisant ; c) $\operatorname{Ext}^n(-, M)$
exact et $0 < \mathrm{lg}_{\mathfrak{p}} M \leqslant \mathrm{lg}_{\mathfrak{p}}
A$ ; b') $H^n(M) \simeq I$ ; c') $H^n(M)$ injectif et la même condition de
longueur ; d) $\operatorname{Ext}^{n+1}(k, M) = 0$ et la même condition ;
\item[(i bis)] $\operatorname{Ext}^{n-1}(k, M) = 0$, $\operatorname{Ext}^n(k, M)
\simeq k$ et $M$ fidèle ;
\item[(ii)] $\operatorname{Ext}^n(-, M)$ dualisant, ou exact avec la condition
de longueur ;
\item[(iii)] $\operatorname{Ext}^n(k, M) \simeq k$ et $M$ fidèle ;
\item[(iv)] $\operatorname{Ext}^n(A/\mathfrak{q}, M) \simeq D(A/\mathfrak{q})$
et $M$ fidèle --- « peut-être suffit-il ici que [ce module] soit comonogène et
$M$ de support $\operatorname{Spec}(A)$ » ;
\item[(v)] le socle de $M/(x_1M + \cdots + x_nM)$ est de dimension $1$, et $M$
est fidèle\footnote{La page écrit « l'annulateur de $\mathfrak{N}$ », où
$\mathfrak{N}$ est biffé, puis « l'annulateur de $\mathfrak{m}$ dans … est de
dim 1 » : c'est le socle, de dimension $1$ sur $k$. Une variante f) de (i bis)
est barrée ; des doubles flèches à gauche relient (i) à (ii) et à (iv).}.
\end{enumerate}
Toutes sont vraies si $A$ est de Cohen-Macaulay et $M \simeq \Omega_A$. La
page esquisse la réciproque et démontre les équivalences internes à (i) --- sous
la réserve faite plus haut sur $\operatorname{Ext}^{n-1}$, c'est-à-dire pour $M$
de Cohen-Macaulay de dimension $n$ : b $\Rightarrow$ b', c $\Rightarrow$ c',
a $\Rightarrow$ b sont triviaux ; c' $\Rightarrow$ b' parce que $D(H^n(M))
\simeq \widehat{A}^m$ et que la condition de longueur force $m = 1$ ; c
$\Rightarrow$ d par l'exactitude à gauche de $\operatorname{Ext}^{n+1}(-, M)$ ;
d $\Rightarrow$ c' ; b' $\Rightarrow$ $M$ fidèle, puisque $H^n(M) = I$ l'est.
Enfin, pour $M$ de Cohen-Macaulay, $\operatorname{Ext}^n(k, M)$ est le socle de
$\operatorname{Ext}^n(A/\mathfrak{q}, M)$, d'où f) $\Rightarrow$ e), et e)
donne $H^n(M) \subset I$ ; l'annulateur de $H^n(M)$ est celui de $M$, $M$
étant de Cohen-Macaulay, et $M$ étant fidèle il est nul, de sorte que, $A$
supposé complet, $H^n(M) = I$\footnote{La page écrit
$\operatorname{Ext}^n(A/\mathfrak{q}, M) = \operatorname{Hom}(A/\mathfrak{q},
M)$ et $\operatorname{Ext}^n(k, M) = \operatorname{Hom}(k, M)$, où l'on attend
$M/(x_1, \ldots, x_n)M$ au second argument ; un passage intermédiaire est
encadré et biffé.}.

\pagerange{36}{37}
\subsection*{Les problèmes (pages 36 et 37)}

\textbf{Problème A.} Si le socle de $M/(x_1, \ldots, x_n)M$ est de dimension
$1$ et $\dim M = n$, montrer que le socle de $M$ est nul.

\textbf{Problème B.} Si le socle de $\operatorname{Ext}^n(A/\mathfrak{q}, M)$
est de dimension $1$ et $\dim M = \dim A$, montrer que le socle de $M$ et celui
de $A$ sont nuls.

Il en tirerait que $M$ est de Cohen-Macaulay et que, $A'$ désignant l'anneau
des homothéties de $M$, $A'$ est de Cohen-Macaulay et $M \simeq \Omega_{A'}$.
Le problème A a une réponse négative : pour $A = M = k[[x,y]]/(x^2, xy)$, de
dimension $1$, et le paramètre $y$, $M/yM = k[x]/(x^2)$ a un socle de dimension
$1$, et $x$ est un élément non nul du socle de $M$\footnote{L'exemple est de
cette lecture. Le théorème de la page 38 ajoute, pour (v), l'hypothèse « $M$
de Cohen-Macaulay », qui rend la question sans objet. Nous ne tranchons pas le
problème B.}.

\textbf{Problème C.} Supposons $\operatorname{Ext}^{n-1}(k, M) =
\operatorname{Ext}^{n+1}(k, M) = 0$ et $\operatorname{Ext}^n(k, M) \neq 0$
(c'est-à-dire $H^n(M) \neq 0$). Est-il vrai que $A$ est de Cohen-Macaulay et
$M \simeq \Omega_A^m$ ? Il y a un homomorphisme naturel $M \to \Omega_A^m$, et
il faut montrer que c'est un isomorphisme ; on peut aussi, par dévissage, se
ramener à montrer que les socles de $M$ et de $A$ sont nuls.

Pour $M$ de type fini, la réponse est oui, avec $\widehat{M} \simeq
\Omega_{\widehat{A}}^m$\footnote{Par la note de la proposition 13,
$\operatorname{prof} M \geqslant n$ ; la nullité de $\mu_{n+1}(M)$ entraîne
alors que la dimension injective de $M$ est finie, égale à $n$ puisque
$\mu_n(M) \neq 0$, donc à $\operatorname{prof} A$, et l'on conclut comme pour
la proposition 14. Ces résultats (Bass 1963 ; Fossum, Foxby, Griffith, Reiten
1975 ; Roberts) sont postérieurs aux pages et cités de mémoire.}.

\pagerange{38}{39}
\subsection*{Le théorème (pages 38 et 39)}

\textbf{Théorème.} Les implications marquées sur le diagramme des conditions
(page 35) sont vraies. De plus, si dans (i) et (i bis) on remplace la condition
$\operatorname{Ext}^{n-1}(k, M) = 0$ par « $M$ est de Cohen-Macaulay », et si
dans (v) on suppose de plus $M$ de Cohen-Macaulay, alors (i), (i bis) et (v)
deviennent équivalentes à : \emph{$A$ est de Cohen-Macaulay et $\widehat{M}
\simeq \Omega_{\widehat{A}}$} --- et de même (ii), (iii), (iv), qui pour $M$
de Cohen-Macaulay en sont des variantes\footnote{Le début de l'énoncé est en
partie illisible ; nous écrivons $\widehat{M} \simeq \Omega_{\widehat{A}}$ là
où la page écrit $M \simeq \Omega_A$, la démonstration supposant $A$ complet.}.

\emph{Démonstration.} $M$ est équidimensionnel, donc (i) $\Leftrightarrow$
(i bis). Récurrence sur $n$ ; pour $n = 0$, $H^0(M) = M \simeq I$, et c'est
trivial. Pour $n > 0$ : $H^n(M)$ est fidèle, donc $M$ l'est ; $A$ se plonge
dans une somme finie de copies de $M$, de sorte que les premiers associés de
$A$ sont parmi ceux de $M$, tous de dimension $n$ puisque $M$ est de
Cohen-Macaulay. Il y a donc un $x \in \mathfrak{m}$ hors de tous, non diviseur
de zéro dans $A$ ni dans $M$. Le lemme de Rees (page 29) donne
$\operatorname{Ext}^i_A(N, M) \simeq \operatorname{Ext}^{i-1}_{A/xA}(N, M/xM)$
pour tout $A/xA$-module $N$, donc la condition (i a) passe à $M/xM$ sur $A/xA$,
qui est encore de Cohen-Macaulay, et la récurrence s'applique : $A/xA$ est de
Cohen-Macaulay, donc $A$ aussi, et $M/xM$ est fondamental pour $A/xA$. D'où
\[
  D(M/xM) = \operatorname{Ker}(x \mid D(M)) \simeq H^{n-1}(A/xA) =
  \operatorname{Ker}(x \mid H^n(A))
\]
(proposition 10)\footnote{La page écrit $H^n(A/xA)$ ; c'est $H^{n-1}$.}.
L'homomorphisme naturel $H^n(A) \to D(M)$, déduit de $H^n(M) \simeq H^n(A)
\otimes M \simeq I$, induit donc un isomorphisme sur les éléments annulés par
$x$ ; comme $H^n(A)$ est de $x$-torsion, il est injectif, et par transposition
$\widehat{M} \to \Omega_{\widehat{A}}$ est surjectif. Il induit un isomorphisme
sur les $H^n$, de sorte que le noyau est de dimension $< n$ ; $M$ étant
équidimensionnel, ce noyau est nul.

Pour (v), $M$ de Cohen-Macaulay et fidèle : $\operatorname{Ext}^n(k, M) \simeq
\operatorname{Hom}(k, M/(x_1, \ldots, x_n)M) \simeq k$, ce qui ramène à (i bis)
--- « cf.\ § 1, prop.\ 4, cor.\ 1 », hors du dossier.

\pagerange{40}{42}
\section*{V. Seconde récapitulation, et le cas $M = A$ (pages 40 à 42)}

\pagerange{40}{40}
\subsection*{Les conditions, reprises (page 40)}

La page 40, écrite à l'encre noire puis reprise à l'encre bleue et au crayon,
fixe $A$ quotient d'un anneau régulier $A'$, et récrit la liste : (i)
$\operatorname{Ext}^n(k, M) \simeq k$, $\operatorname{Ext}^{n+1}(k, M) =
\operatorname{Ext}^{n-1}(k, M) = 0$ ; (ii) $\operatorname{Ext}^n(k, M) \simeq
k$ et $M$ fidèle ; (ii bis) $\operatorname{Ext}^n(A/\mathfrak{q}, M) \simeq
D(A/\mathfrak{q})$ et $M$ fidèle, avec trois questions --- faut-il
$\operatorname{Ext}^{n-1}(k, M) = 0$, ou $A$ équidimensionnel ? suffit-il que
ce module soit comonogène ? ; (iii) $\operatorname{Ext}^{n-1}(k, M) = 0$ et
l'une des conditions équivalentes $\operatorname{Ext}^{n+1}(k, M) = 0$,
$H^n(M) \simeq I$, $H^n(M)$ injectif, avec la condition de longueur ; (iv)
$\operatorname{Ext}^n(-, M)$ dualisant ; (iv bis) exact, avec la condition de
longueur ; (v) « faible » : le socle de $M/\sum x_iM$ est de dimension $1$ et
$M$ est fidèle\footnote{Dans (ii bis), la page écrit
$\widehat{A/\mathfrak{q}}$. Nous lisons le chapeau comme le dual $D$, par
comparaison avec la condition (iv) de la page 35, et de même pages 41 et 42 ;
page 22, en revanche, il désigne le complété. Trois reprises, (iii), (i bis),
(iv), sont encadrées et barrées de zigzags.}. La conclusion visée est encadrée
: « $A$ est C.-M., et $M$ est $\simeq \Omega_A$ ».

Puis le problème, encadré :

\textbf{Pb.} Si $M \neq 0$, $\operatorname{Ext}^{n-1}(k, M) = 0$ et
$\operatorname{Ext}^{n+1}(k, M) = 0$ (ou $H^n(M)$ injectif), $A$ est-il de
Cohen-Macaulay avec $M \simeq \Omega_A^m$ ?

C'est le problème C, sans l'hypothèse $\operatorname{Ext}^n(k, M) \neq 0$ ;
pour $M$ de type fini la réponse est encore oui, $A$ étant ici quotient d'un
anneau régulier\footnote{Si $\operatorname{Ext}^n(k, M) = 0$, la dimension
injective de $M$ serait finie et $\leqslant n - 2$, ce que le théorème
d'intersection de Roberts (1987) interdit pour $M \neq 0$. Cité de mémoire.
Une note en coin, en partie illisible, évoque « $M \simeq \Omega_A \times
A/\mathfrak{I}$ » ; les notes obliques de la marge gauche, en deux couches, ne
se lisent qu'en fragments, et un dernier paragraphe est barré.}.

\pagerange{41}{41}
\subsection*{L'exactitude des foncteurs $\mathrm{Ext}$ (page 41)}

\textbf{Proposition.} Soit $M$ un $A$-module. Si le foncteur $N \mapsto
\operatorname{Ext}^i(N, M)$, sur les modules $N$ de longueur finie, est exact
à gauche --- il change une suite exacte $0 \to N' \to N \to N'' \to 0$ en une
suite exacte $0 \to \operatorname{Ext}^i(N'', M) \to \operatorname{Ext}^i(N, M)
\to \operatorname{Ext}^i(N', M)$ ---, alors
\[
  \operatorname{Ext}^i(N, M) \simeq \operatorname{Hom}(N, H^i(M)).
\]
En effet, un foncteur contravariant exact à gauche sur les modules de longueur
finie est représenté par $\varinjlim_r \operatorname{Ext}^i(A/\mathfrak{m}^r, M)
= H^i(M)$\footnote{L'argument est de cette lecture ; la page n'en donne pas.
« exact à g. » est une lecture incertaine. Une question en marge, en partie
biffée, demande si l'on a alors $\operatorname{Ext}^{i-1}(k, M) = 0$.}.

Si le foncteur est \emph{exact}, $H^i(M)$ est injectif ; pour $M$ de type fini
avec $\dim M \leqslant n$, ou bien il est nul (ce qui arrive pour $i > \dim M$
ou $i < \operatorname{prof} M$), ou bien $i = n$ et $D(H^n(M)) \simeq
\widehat{A}^m$. Alors $M$ est fidèle, $\operatorname{Ext}^n(k, M) \simeq k^m$,
et, $\operatorname{Ext}^{n+1}(-, M)$ étant à son tour exact à gauche, donc nul,
$\operatorname{Ext}^i(k, M) = 0$ pour $i > n$. Une addition à l'encre noire
demande si $A$ est alors de Cohen-Macaulay avec $M \simeq \Omega_A^m$ : c'est
encore le problème C\footnote{La page écrit $\widehat{H^n(M)} \simeq A^m$ ;
voir la note sur le chapeau, page 40. Que $i = n$ vient de ce que
$D(H^i(M))$ est de dimension $\leqslant i$ ; la page ne le justifie pas.}.

Enfin, pour $M$ quelconque, pas nécessairement de type fini : si $N \mapsto
\operatorname{Ext}^{i-1}(N, M)$ est exact à droite, le foncteur
$\operatorname{Ext}^i(-, M)$ est exact à gauche, et donc
$\operatorname{Ext}^i(k, M) = \operatorname{Hom}(k, H^i(M))$ est nul si et
seulement si $H^i(M) = 0$, $H^i(M)$ étant de $\mathfrak{m}$-torsion\footnote{La
page énonce une alternative : « ou bien $\operatorname{Ext}^{i-1}(k, M) = 0$,
ou bien $\operatorname{Ext}^i(k, M) = 0 \Leftrightarrow H^i(M) = 0$ ». La
seconde branche vaut toujours sous l'hypothèse, et l'alternative est inutile.
Les trois exposants de ce paragraphe sont surchargés et douteux.}. La page
s'arrête sur une suite $0 \to M \to J \to M' \to 0$, début vraisemblable d'un
décalage par un injectif.

\pagerange{42}{42}
\subsection*{Le cas $M = A$ (page 42)}

À l'encre gris-noir, la page récrit les conditions pour $M = A$ :
\begin{enumerate}
\item[(i)] $\operatorname{Ext}^n(k, A) \simeq k$ ;
\item[(ii)] $\operatorname{Ext}^{n-1}(k, A) = 0$ ;
\item[(iii)] $\operatorname{Ext}^{n+1}(k, A) = 0$ ;
\item[(iii bis)] $\operatorname{Ext}^n(A/\mathfrak{q}, A) \simeq D(A/\mathfrak{q})$ ;
\item[(iv)] $N \mapsto \operatorname{Ext}^n(N, A)$ est dualisant ; (iv bis) il
est exact ;
\item[(v)] $H^n(A) \simeq I$ ; (v bis) $H^n(A)$ est injectif ; (v ter) $H^n(A)$
est comonogène et $A$ équidimensionnel ;
\item[(vi)] $A/\sum x_iA$ est autodual ;
\item[(vii)] $A \simeq \Omega_A$\footnote{La page écrit
« $\overline{A} \simeq \overline{\Omega}_A$ », et la barre n'est définie nulle
part. Le schéma d'implications donne (vii) équivalente à (v) ; nous lisons
donc (vii) comme « $A$ est fondamental pour lui-même », la condition (a) de la
proposition 15. Une première condition (i bis) est biffée.}.
\end{enumerate}
avec le schéma d'implications
\[
  \bigl[(\mathrm{iv\ bis}) \Leftrightarrow (\mathrm{iv})\bigr] \Rightarrow
  \bigl[(\mathrm{vii}) \Leftrightarrow (\mathrm{v}) \Leftrightarrow
  (\mathrm{v\ bis}) \Leftrightarrow (\mathrm{v\ ter})\bigr] \Rightarrow
  (\mathrm{iii}),
  \qquad
  \bigl[(\mathrm{iv\ bis}) \Leftrightarrow (\mathrm{iv})\bigr] \Rightarrow
  (\mathrm{i}),\ (\mathrm{vi}).
\]
Par la proposition 15, le crochet central est juste à condition de lire (v
ter) comme dans sa condition (d) corrigée. La flèche vers (iii), elle, ne vaut
que si $A$ est de Cohen-Macaulay\footnote{Pour le cône sur une surface
abélienne de la page 33, $A \simeq \Omega_A$ mais $A$ n'est pas de
Gorenstein, sa dimension injective est infinie, et la non-nullité des nombres
de Bass au-delà de la profondeur donne $\operatorname{Ext}^{n+1}(k, A) \neq
0$. Cité de mémoire. Le schéma est redessiné ; un pavé d'encre couvre ce qui
était écrit entre la flèche horizontale et le cadre de droite.}.

Sous (vii) écrit « $A$ est C-M.\ et $\Omega_A = A$ », la page entreprend de
prouver, par récurrence sur $n$, que \emph{(vii) entraîne que $A$ est de
Cohen-Macaulay} --- l'énoncé du théorème 5 biffé de la page 33. $A$ est
équidimensionnel puisque $\Omega_A$ l'est ; soit $x \in \mathfrak{m}$ non
diviseur de zéro dans $A = \Omega_A$. Il suffirait de montrer
\[
  \Omega_{A/xA} \simeq \Omega_A/x\Omega_A ,
\]
car alors $A/xA$ serait fondamental pour lui-même et la récurrence
s'appliquerait. Pour le voir, $A$ étant quotient d'un anneau régulier $A'$, le
lemme 6 appliqué à $A \to A/xA$ et à $\operatorname{Ext}_{A'}(-, A')$ donne une
suite spectrale
\[
  \operatorname{Ext}^i_A\bigl(A/xA, \Omega_A^{(j)}\bigr) \Rightarrow
  \Omega^{(\ast)}_{A/xA},
\]
dont seules les colonnes $i = 0, 1$ sont non nulles ; d'où la suite exacte
sur laquelle le dossier s'arrête :
\[
  0 \to \operatorname{Ext}^1_A(A/xA, \Omega_A) \to \Omega_{A/xA} \to
  \operatorname{Hom}_A(A/xA, \Omega_A^{(1)}) \to 0 .
\]
Le premier terme est $\Omega_A/x\Omega_A$. L'égalité voulue demande donc que
$x$ ne soit pas diviseur de zéro dans $\Omega_A^{(1)}$, le module de déficience
dual de $H^{n-1}(A)$ (page 28) --- et c'est exactement ce qui peut manquer.
Pour le cône sur une surface abélienne, $n = 3$, $H^2(A)$ est de longueur
finie et non nul, $\Omega_A^{(1)}$ aussi, et aucun $x \in \mathfrak{m}$ n'y est
non diviseur de zéro : la démonstration ne peut pas aboutir, et l'énoncé est
faux\footnote{L'analyse de l'obstruction et l'exemple sont de cette lecture.
Sur la page, l'exposant de $\operatorname{Ext}_{A'}$ est incertain (« $r - n$ »,
« $r - n + 1$ » au-dessus) ; $\Omega_{A/x}$ est souligné d'un trait incurvé, et
un crochet en marge droite, avec la note « Lemme », ne se lit qu'en partie. Le
dossier s'arrête sur cette ligne.}. Avec l'hypothèse que $A$ est de
Cohen-Macaulay, $\Omega_A^{(1)} = 0$, la suite donne l'isomorphisme, et l'on
retrouve la proposition 10 : c'est la définition 2 de la page 34, qui exige
Cohen-Macaulay, qui avait raison.

\end{document}
